% 原书第 439 页（PDF 第 452 页）
\BookPage{439}{452}
\BookHeading{康雍干时期台湾的乡治}{康雍干时期台湾的乡治}{杨晋平*}
\BookSubheading{一、前言}

清史学者萧一山在《清代通史》中，对清代的乡治有一段精辟论述，他说：

\begin{BookQuote}
盖清代之地方制度，自县以下，委诸民治，官书所载，殊不甚详，或里或乡，或保或社，或区或村，一地迭用数种名称者有之，数种名称而各地互异者亦有之，论其组织之系统如何？行政方式如何，要皆莫能详解，初惟就习惯上之方便，随意采用，固未尝依合理之规制，具整齐划一之方式。里乡保社，应用之名称既杂，故后之论者，乃有一体二体，或一体两端之不同。1
\end{BookQuote}

清代行政层级只到厅、县，虽有县丞、巡检、典史等属官，但一个县，辖境辽阔，以台湾为例，编制较大的嘉义县官员、吏员、差役共二百二十一人，而辖境从台中大甲到基隆三貂几乎占台湾一半版图的淡水厅，编制仅一百九十四人。²在人力与组织两缺的情况下，「委诸民治」是必然的，所以在清领台湾的二百一十二年间，台湾乡治一定受到倚重，因此有「台属有司，率多不亲理民事，皆批乡约、保长、管事、书役查覆。」3的现象出现。这也说明清代台湾乡治组织必相当发达，也十分重要，但可惜「官书所载，殊不甚详」，历史研究者首重的就是史料，受限于档案缺乏，史料不足，致使清代台湾乡治研究不容易开展，研究者不多，研究成果不丰。

清代台湾县级以下的地方制度「其组织系统如何？行政方式如何？」至今还是「莫能详解」，如能试图去解，这必是浩大工作，不是本文有能力处理，现只选择

\vfill
\begin{BookNotes}
\BookNote{*}{佛光大学历史所硕士。}
\BookNote{1.}{萧一山，《清代通史·第四篇》（台北：商务印书馆，1976年），页629。}
\BookNote{2.}{张胜彦，《清代台湾厅、县制度之研究》（台北：华世出版社，1993年），页2、74、98。}
\BookNote{3.}{高其倬等撰，《清经世文编选录》(南投:台湾省文献委员会，1979年)，页66。}
\end{BookNotes}

% 原书第 440 页（PDF 第 453 页）
\BookPage{440}{453}
\noindent 《康雍干时期台湾的乡治》做为课题，希望以此切入，能有所得，也希能抛砖引玉，引起更多学子注意，解开萧一山先生「莫能详解」的遗憾，让清代初期台湾的乡治能有一较清晰、完整的轮廓。

\BookSubheading{二、清初台湾的治理}

台湾在中国的边陲，又隔大海雍正就说:「台湾地方，自古未属中国；皇考圣略神威，取入版图」。4清是在康熙二十二年（1683）八月，派福建水师提督施琅攻打台湾，郑克坏出降，才开始领有台湾。自康熙（1683）经雍正到乾隆（1795）三朝，一百一十二年，以当时的交通与通讯台湾与中国相距遥远，台湾受到清朝的关注不多，遑论重视，就如康熙所言：「台湾属海外地方，无甚关系；因从未向化，肆行骚扰，滨海居民迄无宁日，故兴师进勦。即台湾未顺，亦不足为治道之缺」。并认为「台湾仅弹丸之地，得之无所加、不得无所损」。所以清初对台湾的治理，注重在消极的海防安全，并无积极经营意图，一切措施政策皆仅基于治安上的考虑，是属闭锁管控格局。所以有台湾去留的议论，厅、县设立的被动及消极治台的作为。

\BookSubheading{（一）台湾的去留}

当清收有台湾，中央与地方官员对台湾地位和海防的重要性，缺乏真正的了解，「台湾应弃应守」曾经过一段时间的议论，朝廷和地方都、抚多主张放弃台湾，所以有「台湾初定，提督施琅请设官镇守，廷议未决，有谓宜迁其人、弃其地者」，8以为郑氏既降「恃为窟穴，倚险负固，飘突靡常」就不会再有，更有以「台湾削此一块荒壤，无用之地耳，去之可也。」10等言论，当时的康熙也是举棋不定，后经施

\vfill
\begin{BookNotes}
\BookNote{4.}{《清世宗实录选辑》（南投：台湾省文献委员会，1997年)，页4。}
\BookNote{5.}{《清圣祖实录选辑》（南投:台湾省文献委员会，1997年)，页129。}
\BookNote{6.}{同上注，页129。}
\BookNote{7.}{同上注，页127。}
\BookNote{8.}{《清史稿台湾资料集辑》（南投:台湾省文献委员会，1994年)，页461。}
\BookNote{9.}{《清代官书记明台湾郑氏亡事》（南投:台湾省文献委员会，1995年），页29。}
\BookNote{10.}{施琅，《靖海纪事》（南投:台湾省文献委员会，1995年），页62。}
\end{BookNotes}

% 原书第 441 页（PDF 第 454 页）
\BookPage{441}{454}
\noindent 琅以「弃之必为外国所踞；奸宄之徒窜匿其中」11力争，乃在台湾设一府三县及巡道一员。并置总兵官一员、副将二员、兵八千，分为水陆八营。另澎湖设副将一员、兵三千，分为二营，12屯兵驻防。

\BookSubheading{（二）府县的设置}

汉人对台湾大规模的拓垦，开始于1624年荷兰入据台湾后，因荷兰人实行王田政策，大量招汉佃来台开垦，而开垦以安平为中心，以台南周边北港、萧垅、湾里、茄拔、新港、大目降为范围，最远到高雄冈山的阿公店。清朝领台初期仍以台南为中心，并依明郑旧制将承天府改置台湾府，并设台湾县、另将北路天兴州改为诸罗县，南路万年州改为凤山县。其治理范围与规制没有超越明郑时期。当时疆域是以府治为中心，东至保大里大脚山五十里为界，西至海为界。南至沙马矶头六百三十里为界，是南路；矶以内诸社。北至鸡笼山二千三百一十五里为界，是北路；土番居多。惟近府治，是汉、番参半。至于东方，山外青山，迤南亘北，皆不奉教。生番出没其中，人迹不经之地；延袤广狭，莫可测识13。

在康熙后期以后，南部可耕地逐渐被垦透，汉人垦民逐渐北移。垦首施世榜、杨申、吴洛、黄仕卿等又集资开凿八宝圳、十五庄圳、恶马圳，使垦务更为发展，康熙晚期垦成二水、田中、社头、大村、花坛等地。雍正年间拓成的有芬园、北斗、鹿港。乾隆年间拓垦再往北发展，以员林、埔心、永靖、埔盐、溪湖为重心，并向台中拓垦头店·大墩发展。

雍正元年(1723)，为因应诸罗县以北纷至踏来的垦民，民番杂处等问题，请增设知县一员、典史一员；并请增设捕盗同知一员，14朝廷探纳后，乃析诸罗县以北，半线地方另设彰化一县，并增淡水捕盗同知一员；雍正九年（1731），又将彰化县大甲溪以北钱粮、命盗等项归淡水同知就近管理，到雍正十一年（1733）并移淡水同知

\vfill
\begin{BookNotes}
\BookNote{11.}{《清圣祖实录选辑》（南投:台湾省文献委员会，1997年），页131。}
\BookNote{12.}{同上注，页132。}
\BookNote{13.}{高拱干，《台湾府志》（台北:台湾银行经济研究室，1960年)，页6。}
\BookNote{14.}{《清世宗实录选辑》（南投:台湾省文献委员会，1997年），页3-4。}
\end{BookNotes}

% 原书第 442 页（PDF 第 455 页）
\BookPage{442}{455}
\noindent 驻竹堑。

康雍干三朝在台湾的治理仅限在西部沿海一带及浅山地区，东北部的宜兰要晚至嘉庆后期，山区的南投要到道光年间，后山的花莲、台东设治则是在清末的事。

\BookSubheading{（三）消极治理}

清初在台治理版图是呈一种扩张趋势。这扩张是由南部向北部逐渐拓展，由沿海向内山推进；清政府祇是随人民的自然发展，被动的设官治理而已。政府并无扩张意图，反作种种局限措施。例如康雍二朝严格限制汉人移民来台，到乾隆朝，对于内地人民的渡台，犹时而限制、时而放宽。1为治安考量，所颁布的渡台禁令，限制汉人渡台，延迟了台湾开发与发展。对入台者，也划定拓垦的范围与区域，以土牛为界，严禁越界私垦，也因此台湾拓垦是在阻力重重，违法犯禁下艰困进行。另清朝为控制上的方便，不许台湾府县筑成，置各级官员与民众安全于不顾。又台湾官吏不良，朝廷也深知，但长期放纵，致使官逼民反，三年一小乱，五年一大乱。

1、渡台禁令

清领有台湾后对汉人移民是探取限制、干扰政策。荷兰与明郑时期奖励垦殖的措施，被弃置不用，反而「克台之岁，旨下福建督抚，凡渡台者禁带家眷，而琅亦请申海禁，不许惠、潮之人入台」，16有学者将这些政策归纳成三点，称「移民三禁」也有简称「渡台禁令」17；其主要内容是：一是惠、潮是海盗渊薮，禁止粤人渡台。二是渡台必需要官府「印单」，出入严格检查，禁止自由来台。三是渡台只限单身，已渡台者不得招来家眷18。如此禁制政策，使清代台湾偷渡非常盛行。

康熙时，为限制汉人来台，由该县给发照单，如地方官滥给来台湾照单，经督、

\vfill
\begin{BookNotes}
\BookNote{15.}{《清高宗实录选辑》（南投:台湾省文献委员会，1997年），页1。}
\BookNote{16.}{连横，《台湾通史》（南投:台湾省文献委员会，1992年），页152。}
\BookNote{17.}{清朝厉行渡台禁令政策时间相当长，其间仅有短暂的数次弛禁。大致说来，清朝统治台湾的211年当中，从1684年到1790年之间是采取一种较严格的禁止与限制；1790年以后，才较放松，到了1875年因沈葆桢建议，才真正开放移民。}
\BookNote{18.}{郭廷以，《台湾史事概说》（台北:正中书局，1996年），页124。}
\end{BookNotes}

% 原书第 443 页（PDF 第 456 页）
\BookPage{443}{456}
\noindent 抚题参一次者，罚俸六月；二次者，罚俸一年；三次者，降一级留任；四次者，降一级调用。雍正为严禁偷渡来台，对失察水汛及本地文武各官，罚俸一年，如文武衙门隐匿不报者，或被告发、或被上司查参，将该管兼专文武各官，皆照讳盗例分别议处。19对船户管制尤为严厉，如有违禁揽载偷渡者，一经拏获，即严行究拟。兵役致疏脱十名以上者，兵役各责二十板；每过十名，加责十板；至四十名以上，各责四十板，革役；五十名以上，枷号一月发落。而在回船时，如有照外多载或顶充偷渡、又潜匿报，船户、舵水俱照窝藏盗贼例治罪；出结之族邻、行保，杖一百、徒三年。20

海禁渐严，在康熙五十七年（1718）时，凡往来台湾之人，必令地方官给照，方许渡载；单身游民无照者，不许偷渡。21致使在台立业者，既不能弃其田园，又不能搬移眷属；官方恐其等终岁群居，皆无家室，则其心不靖，难以久安，在雍正十年(1732）填给路引，准许搬眷入台。到乾隆五年（1740)以流寓民眷均已搬取为由，停止给照，不准搬移。乾隆十一年（1746）又准祖父母、父母及妻子欲赴台侍奉就养者，仍准其给照搬养，次年又不准给照，凡有渡台民人，禁绝往来，不能搬移。22

乾隆二十五年（1760）又议给照，搬眷。乾隆三十四年（1769）暂署福州将军崔应阶因黄教民变上奏「台湾流寓内，闽人约数十万、粤人约十余万，而渡台者仍源源不绝。此皆穷困逋逃之辈，性情狡悍，不能安分；结伙连群，势必滋生事端。与其办理于临时，不如羁縻于平日；现在严加稽察，有罪者酌量遣发，并设法钤束」等语。乾隆则谕以：

\begin{BookQuote}
所见固是；而于正本清源之道，犹未尽合。此等渡台民人，多属内地素无恒产、游手好闲之徒；一经潜渡海洋、窜迹台地，日积日多，必致引类呼朋，毫无顾忌。黄教之案，乃其明验。但此等无赖游民与其约束于到台之后，多费周章；不如稽察于渡海之前，力为禁阻。向来腹地民人不许私行赴台，定例本严；因地方有司奏行不力，以致诡名偷越之人日增一日。否则，此十余万及数
\end{BookQuote}

\vfill
\begin{BookNotes}
\BookNote{19.}{《清会典台湾事例》(台北:台湾银行经济研究室，1966年)，页30。}
\BookNote{20.}{同上注，页168-169。}
\BookNote{21.}{《清圣祖实录选辑》(南投:台湾省文献委员会，1997年)，页168。}
\BookNote{22.}{余文仪，《续修台湾府志》（南投：台湾省文献委员会，1993年），页725-726}
\end{BookNotes}

% 原书第 444 页（PDF 第 457 页）
\BookPage{444}{457}
\begin{BookQuote}
十万闽、粤流寓、岂能不胫而至耶。23
\end{BookQuote}

乾隆四十九年（1784）上谕复以「渡台民人，多属内地游手好闲之徒，一经窜迹台地，日积日多，必致引类呼朋，毫无顾忌。着将福建、广东赴台人民，严行禁止，设法巡逻，毋许私行透漏。

在《清会典事例》对渡台民人也有严厉管制规定：

\begin{BookQuote}
台湾流寓之民凡无妻室者，应逐令过水，交原籍收管。其有妻子、田产者，如犯歃血订盟、诱番杀人、捏造匿名揭帖、强盗窝家、造卖赌具，应拟斩、绞、军、流等条，除本犯依律例定拟外，此内为从罪轻之人并教唆之讼师，均应审明逐令过水。其越界生事之汉奸，如在生番地方谋占番田并句串棍徒包揽偷渡及贩卖鸦片者，亦分别治罪，逐令过水。24
\end{BookQuote}

渡台禁令时禁时弛，开禁仅准人民搬眷，而非准许人民自由来往。在台有业良民欲自内地搬眷，或内地人民欲至台湾探亲，均须经过麻烦的领照手续。只身无业之民，在台无亲属可依者均不准渡台。

在禁令实施的间，限制了大陆与台湾人民往来自由，不少人被迫铤而走险，即行偷渡来台，造成两地海防的严重漏洞。民人为辗转偷渡来台，甘冒生死风险，其间，种种人间惨剧接连发生，并沿生偷渡有关暗语「灌水」、「放生」、「种芋」、「饵鱼」，足见当时偷渡盛行。25

清初的渡台禁令，是以谋求台湾的安定为首要目标，惟就其政策设计、渡台禁令的实施及影响而言，显然并未达成想定目标，只因当时「漳、泉内地无籍之民无田可耕、无工可雇、无食可觅，一到台地，上之可以致富，下之可以温饱，一切农工商贾以及百艺之末，计工授直，比内地率皆倍蓰。而必曰尔其坚坐饿死，无往求生为也」。26所以渡台势不可禁，清朝未能体察人民所苦，强势禁止，反引起如政治、治安、社会等诸问题，足证其治台政策确有可再斟酌的地方。

\vfill
\begin{BookNotes}
\BookNote{23.}{《清高宗实录选辑》（南投:台湾省文献委员会，1997年），页198。}
\BookNote{24.}{《清会典台湾事例》（台北:台湾银行经济研究室，1966年），页168。}
\BookNote{25.}{王瑛曾，《重修凤山县志》（南投:台湾省文献委员会，1993年），页212。}
\BookNote{26.}{高其倬等，《清经世文编选录》（南投：台湾省文献委员会，1997年），页2。}
\end{BookNotes}

% 原书第 445 页（PDF 第 458 页）
\BookPage{445}{458}
所以早在雍正时台湾知府沈起元救说：

\begin{BookQuote}
按台郡之入版图，正如洪荒初辟，章程制度，草草未立。而六十年以前官斯土者，因循废弛，既所不免；至朱一贵变后，震动当局，地方大臣未敢以致乱之由据实陈告，乃徒为因噎废食之计，禁官府之携家、禁民户之搬眷、禁内民之渡台、汰道标之守兵，以为是可以久安长治矣，何其远欤？27
\end{BookQuote}

清初的渡台禁令是一错误政策，直到光绪元年（1875），经钦差大臣沈葆桢奏请，施行近二百年的渡台禁令与私入番境同时弛禁，不再禁止。28

2、筑土牛，禁止越界私垦

为限制汉人四出拓垦，原就禁止民人越界私垦，朱一贵民变后蓝鼎元建议：

\begin{BookQuote}
台、凤、诸三县山中居民，尽行驱逐，房舍尽行拆毁，各山口俱用巨木塞断，不许一人出入。山外以十里为界，凡附山十里内民家，俱令迁移他处；田地俱置荒芜。自北路起，至南路止，筑土墙高五、六尺，深挖濠堑，永为定界。越界者以盗贼论。如此则奸民无窝顿之处，而野番不能出为害矣。」29
\end{BookQuote}

雍正三年（1725）福建巡抚臣毛文铨又以「夫生番一种，向不出外，皆潜处于伊界之中，耕耘度活。内地人民，或因开垦而占其空地间山，或因砍木而攘其藤梢、竹木，生番见之，未有不即行杀害、酿成大案者」。30为防止生事，惟有清其域限，严禁诸色人等，不许辄入生番界内，方得无事。所以在逼近生番交界之间，各立大碑，杜其擅入。

乾隆十年（1745）福建布政使司布政使高山上呈〈陈台湾事宜疏〉说：

\begin{BookQuote}
番民地界之宜照旧划清也。查台郡西临大海、东逼崇山，由山至海不过数十里；自南至北，绵互二千余里。崇山之内，皆生番所居；界外平埔，系熟番、汉民零星散处。从前定有地界，立石开沟；久而失址，甚有拔石塡沟，那移改
\end{BookQuote}

\vfill
\begin{BookNotes}
\BookNote{27.}{同上注，页6-7。}
\BookNote{28.}{《台湾中部碑文集成》（南投:台湾省文献委员会，1994年），页114。}
\BookNote{29.}{蓝鼎元，《东征集》（南投:台湾省文献委员会，1997年)，页40。}
\BookNote{30.}{郑喜夫辑，《雍正朱批奏折选辑》（南投:台湾省文献委员会，1997年)，页4。}
\end{BookNotes}

% 原书第 446 页（PDF 第 459 页）
\BookPage{446}{459}
\begin{BookQuote}
徙，希图越界私垦，致启生番戕杀之机。臣思生番群聚内山、熟番错居社地、汉民散处各庄，自应划清界限，各管各地；不得混行出入，相寻衅端。臣于清查庄产之便，顺途勘验。觉从前所立之界，并无山河依阻，易于无凭；但日久相延，又未纷更改易。相应饬令地方官遴委佐杂，眼同各番土目指出现在管理分界之处，再行立表，划清界限；使生番在内、汉民在外，熟番间隔于其中。清界而后，汉民毋许深入山根，生番毋许擅出埔地；则彼此屏迹，断绝往来，自不致生衅滋事矣。31
\end{BookQuote}

乾隆二十三年（1758）台湾道杨景素又以「界内之民侵展禁限，熟番地日促而汉民地渐近。生番既与熟番搆衅，生番亦乘间焚杀汉民。公案界，掘深沟，筑土牛以为之限」。32同年闽浙总督杨应琚也建议：「台民垦种，侵越熟番地界；应查明挑沟，画清界限」。33

乾隆二十五年（1760）闽浙总督杨廷璋条奏：

\begin{BookQuote}
台郡彰化县沿山番界，年来侵垦渐近内地，生番逸出为害。今据该镇、道勘明，于车路旱沟之外，各有溪沟、水圳及外山山根，堪以久远划界。其与溪圳不相接处，挑挖深沟，堆筑土牛为界。至淡防厅一带，从前原定火燄山等界，仅于生番出没之隘口立石为表，余亦未经划清。今酌量地处险要，即以山溪为界；其无山溪处，亦一律挑沟堆土，以分界限。34
\end{BookQuote}

乾隆五十二年（1787）林爽文民变时乾隆还关心「台湾地方，前闻杨景素在道员任内设立土牛，分别内外地界；其是否办理妥协？并着常青会同黄仕简于勦捕事竣后亲往履勘，斟酌妥善，归于善后事宜内一并妥办」。35

乾隆五十三年（1788）又谕令福康安于勦捕林爽文事竣后务将台湾一切善后事宜全行办竣，即在台湾多住一、二年亦所应得。「如内外山地界从前设立土牛，未为

\vfill
\begin{BookNotes}
\BookNote{31.}{《清奏疏选汇》（南投:台湾省文献委员会，1997年），页40-41。}
\BookNote{32.}{李元度，《碑传选集》(南投:台湾省文献委员会，1994年)，页527。}
\BookNote{33.}{《清高宗实录选辑》（南投:台湾省文献委员会，1997年)，页118。}
\BookNote{34.}{同上注，页126。}
\BookNote{35.}{同上注，页329。}
\end{BookNotes}

% 原书第 447 页（PDF 第 460 页）
\BookPage{447}{460}
\noindent 周密；应重加勘定，务令界限分明，勿使日久偷越，以致滋生事端」。36乾隆也认为「台湾疆土既开，民安耕凿，处处皆成膏腴之地；自杨景素议立界限之后，界外良田美产转畀生番，而生番以射生为业、不事耕种，势必内地民人仍往偷垦，日久徒滋事端。」37

康雍干三朝为防范民人越界偷垦，滋事端，一再勘定，务令界限分明，并订有「凡民人偷越定界私入台湾番境者，杖一百；如近番处所偷越深山抽籐、钓鹿、伐木、探等项，杖一百、徒三年。其本管头目钤束不严杖八十，乡保、社长各减一等。巡查不力之直日兵役，杖一百；如有贿纵，计赃从重论」38的重罚。但民人并没有为土牛沟39所限，也没有为不许垦民越界偷垦的法律所吓止，这在乾隆五十三年（1788）福康安在给乾隆皇帝的奏折中，可清楚的看到，福康安说：

\begin{BookQuote}
设立土牛，禁止奸民越界占垦，免滋事端。乃因生聚日繁，民人私向生、熟番黎佃地耕种，价值稍轻者，谓之租；价值稍重者，谓之典卖。熟番等归化日久，渐谙耕作，祗因业经典卖与民，无由取赎。是以各处番地，不特嘉义以南多有侵越，即淡水等处续定土牛之界，亦成虚设。臣福康安追勦贼匪时，周历全郡，所过近山地方，良田弥望，村落相联，多在舆图定界之外40。
\end{BookQuote}

在康熙六十一年（1721）初定土牛沟，乾隆朝也多次续定，官府为求重视并在档案图册以红笔在舆图上划线，以确定番汉界址，但到乾隆乾隆五十三年（1788）土牛已成虚设，土牛线外已是良田弥望，村落相联。所以在台湾的拓垦上，偷渡、越界私垦等非官府可完全掌控的活动非常猖獗，而在官府势力达不到的地方，还是需要组织与纪律，为垦民排难解纷，维持秩序，地方上在官府未设前，就有一定自治组织，进

\vfill
\begin{BookNotes}
\BookNote{36.}{同上注，页544。}
\BookNote{37.}{同上注，页401。}
\BookNote{38.}{《清会典台湾事例》(台北:台湾银行经济研究室，1966年)，页169。}
\BookNote{39.}{「土牛沟」又称「土牛红线」、「土牛」，是始于康熙六十一年（1721），当时清廷统治台湾初期，汉人一批批来到台湾拓垦，汉人常强行越界开垦原住民的山区祖产，引起原住民更多的「出草」行为，双方常有纠争、死伤无数，故原住民立石为先，禁止汉人逾越偷垦。之后，清廷为防止汉民窝藏「番地」，侵垦原住民的土地，也为使「生番」不能逸出为害汉人，乃在山区挖沟筑立一条『长2丈，底阔1丈，高8尺，顶阔6尺；每沟长15丈，阔1丈2尺，深6尺；永禁民人逾越私垦』的「番汉界」，此一山沟所挖出的黄土堆，一堆堆的土丘远看像牛背，当时俗称为「土牛」。}
\BookNote{40.}{清高宗敕撰，《钦定平定台湾纪略》（台北:台湾银行经济研究室，1961年），页989-990。}
\end{BookNotes}

% 原书第 448 页（PDF 第 461 页）
\BookPage{448}{461}
\noindent 行自我管理，因此清代台湾乡治的蓬勃，是有社会环境需要。

3、不建城墙

\begin{BookQuote}
「保障称贤，崇墉比峻；城池之重，古人咏金汤也」，41所以「城池之设，上以
\end{BookQuote}

卫国、下以卫民」，42台湾「郡县既设，城池其首务焉·高垒深沟，以为不虞之备，所谓设险以守其国者也」。43地方也以「今虽民居稠密，市肆纷错，盛平之时，高枕可以无忧，万一有警；将何所恃以为守御之具？所当亟为为筹划，以为久安长治之计者也」，44请求建城，知府高拱干也以：

\begin{BookQuote}
今虽分设郡县，未有城垣捍卫，以备御侮前道周副使曾经详请前任督宪具题建造，然大司空之议覆，只令捐输营葺·工费浩大，非区区台地文武所能割俸竣事，是以竟成道旁筑舍。伏念海外要地，上费皇上几许金钱兵力，方入版图，当内外关重之处，立万年不拔之基，际此升平固无他虑，而设险守固，先赖金汤；牖户绸缪，事非迂阔。相应详请再赐入告，将台湾筑城工费，或准本省生俊捐纳、或于省会别开事例，庶众志成城，捍御有赖。其所关海国封疆，非细
\end{BookQuote}

故也。45

虽经地方、官员为台湾筑城请命，只是朝廷并不为所动。地方只有围木栅，环刺竹，建土城以求自保，所以有「台湾府无城」，46雍正元年，台湾县知县周钟瑄创建，以木栅为城。47台湾县「旧无城，雍正元年，知县周钟瑄创建木栅」。48凤山县在康熙六十一年（1722），筑土城，雍正十二年（1734），环植刺竹为城。诸罗县在康熙四十三年（1704），始建木栅，雍正元年（1723），改建土城。彰化县在雍正

\vfill
\begin{BookNotes}
\BookNote{41.}{刘良璧，《重修福建台湾府志》（台北：台湾银行经济研究室，1961年)，页75。}
\BookNote{42.}{沈茂荫，《苗栗县志》（台北:台湾银行经济研究室，1960年），页33。}
\BookNote{43.}{陈文达，《台湾县志》（台北：台湾银行经济研究室，1961年)，页69。}
\BookNote{44.}{同上注，页69-70。}
\BookNote{45.}{范咸，《重修台湾府志》（台北:台湾银行经济研究室，1961年)，页306。}
\BookNote{46.}{同上注，页61。}
\BookNote{47.}{同上注，页57。}
\BookNote{48.}{谢金銮、郑兼才，《续修台湾县志》（台北：台湾银行经济研究室，1962年），页6。}
\end{BookNotes}

% 原书第 449 页（PDF 第 462 页）
\BookPage{449}{462}
\noindent 十二年（1734），环植刺竹。淡水厅在雍正十一年（1733），环植刺竹。49

在康雍干的坚持下，台湾筑城努力了一百年后，才得以实现。当局不允许台湾建城，有制度初备，百废方兴，未遑议及与建材砖石之属，皆取资于内地，经费无所出等理由。50但这是表面推诿的官式说法，台湾不能筑城实是朝廷以「台贼多自内生，鲜由外至」，51认为台湾兵备，在防范内忧多余外患。康熙六十年（1721），朱一贵起事抗清，府城失陷，乱事厘平后引起许多有识之士请愿筑城，蓝鼎元就说：「台地民番杂处，狼子野心，顽良参半；建筑城池，确不可易」，52建议筑城，但朝廷仍不为所动。诚如雍正所说：「殊不知城垣之设，所以防外患；如必当建城，虽重费何惜？而台湾变乱，率皆自内生，非御外寇比；不但城可以不建，且建城实有所不可也」。53

经历朱一贵与林爽文二次大规模民变，清朝始体认到无城难于防守，变起既烦兵力，又骇众听闻，为求有备无虞，决定筑城，在乾隆五十二年（1787）谕军机大臣等：

\begin{BookQuote}
台湾地方自康熙年间始隶版图，因远隔海洋，该处府厅各县俱未建立城垣。后经朱一贵聚众滋事平定后，总督满保曾议建城；惟时以地处海外，无城虽难于防守，然失之易、复之亦易，是以未经建立以节繁费。但该处久成郡县，与内地无异；而城围尚相沿用竹木编插，不足以垂久远。且此次林爽文纠众猝起，县城遂被砍破，究由莿竹不能防御所致。朕意与其失之复取既烦我兵力，又骇众听闻，何如有城可守而勿失，更为有备无虞乎！54
\end{BookQuote}

乾隆五十三年（1788）林爽文平定后又谕：『现在台湾府、厅、县，应行改建城垣；德成于工程事务素为熟谙，着即驰驿速赴台湾，会同徐嗣曾将该处应办城工悉心

\vfill
\begin{BookNotes}
\BookNote{49.}{范咸，《重修台湾府志》(台北:台湾银行经济研究室，1961年)，页58。}
\BookNote{50.}{陈文达，《台湾县志》（台北:台湾银行经济研究室，1961年)，页69。}
\BookNote{51.}{蓝鼎元，《东征集》（台北:台湾银行经济研究室，1958年)，页38。}
\BookNote{52.}{同上注，页38。}
\BookNote{53.}{范咸，《重修台湾府志》（台北:台湾银行经济研究室，1961年)，页61-62。}
\BookNote{54.}{《清高宗实录选辑》，（南投:台湾省文献委员会，1997年），页326。}
\end{BookNotes}

% 原书第 450 页（PDF 第 463 页）
\BookPage{450}{463}
\noindent 估勘』，55台湾府、厅、县始开始筑城，惟距康熙二十二年（1683），已过一百零五年。

4、吏治不良

康熙时「承平日久，守土恬熙，绝不以吏治民生为意」，56在经过朱一贵民变后，雍正才开始知道「台湾地方远隔重洋，全在道、府、厅、县各得其人；而该员又须熟悉风土情形，弹心办理，于地方始有裨益」。57乾隆也说:「台湾远隔重洋，最关紧要；道、府、厅、县必须才守兼优之员，方能胜任。朕又闻该省督、抚遇有台湾缺出，不问属员才具是否相宜，多以私人调补。而得缺之员从不以冒险渡海为虑，反视之为利薮」。58又说:「今细思徐访，知逆匪林爽文等起事之由，皆因该地方官平日废弛贪黩，视台湾缺分为利薮，不以冒险渡海为畏途，转以得调美缺为喜」。59而「向来台湾道府缺出，俱由督抚奏调。该督抚等因其地土丰饶，每有瞻徇私情，不问属员才具能否，举请调补，俾得侵渔肥蠹；所调之员不以涉险为虞，转以得调美缺为喜，于吏治大有关系」。60雍正、乾隆都已知台湾吏治不良，仍然坐视「台湾府文职官员，平日并不爱民，但知图利苛索」，61并无积极改善作为。

检视台湾吏治之坏，乾隆三十四年（1769）已革凤山县知县方辅悟亏空库项一万三千两、仓谷三千余石之多。同年诸罗县革职知县陶浚，任内亏空库银一万五千余两；续又查出原揭数外，尚有亏空银一万四千余两、谷二万四千余石。也是当年，又发生台湾县知县曾曰琇等亏空仓库一案。²当时台湾仅有四县，就有三知县犯案。在林爽文民变中被杀官员又有「台湾府知府孙景燧、台防同知刘亨基、原任台湾县知县程峻、署诸罗县事原任台防同知董启埏、署诸罗县唐镒声名俱属狼籍，其贪黩敛怨，

\vfill
\begin{BookNotes}
\BookNote{55.}{同上注，页568。}
\BookNote{56.}{蓝鼎元，《平台纪略》（南投:台湾省文献委员会，1997年），页1。}
\BookNote{57.}{《清世宗实录选辑》（南投:台湾省文献委员会，1997年)，页26。}
\BookNote{58.}{《清高宗实录选辑》（南投：台湾省文献委员会，1997年），页383。}
\BookNote{59.}{同上注，页407。}
\BookNote{60.}{清高宗敕撰，《钦定平定台湾纪略》（南投:台湾省文献委员会，1997年)，页52。}
\BookNote{61.}{《清圣祖实录选辑》（南投:台湾省文献委员会，1997年)，页174。}
\BookNote{62.}{《清高宗实录选辑》（南投:台湾省文献委员会，1997年），页195-198。}
\end{BookNotes}

% 原书第 451 页（PDF 第 464 页）
\BookPage{451}{464}
\noindent 被害实由自取」，63其中「唐镒家本贫苦，委署诸罗县事为时未久，即寄回银二千数百两之多；若非在任时肆意贪婪，从何得此多赀」。64可见台湾吏治一般，所以说：「向来台湾官吏侵渔成习，遇事即思浮冒，以图肥蠹。」65也难怪有「台湾吏治黯无天日」66之说。而乾隆十八年（1753）大浪泵之刘和林树旗，旗内就书「统领八社番民以翦贪官」等字。67后来有人以「各省吏治之坏，至闽而极；闽中吏治之坏，至台湾而极」，68台湾吏治真是不良。

\BookSubheading{三、清初台湾的乡治}

清代县以下乡里组织有「乡、保、里、社、区、总、图、都」69等名称，各地都有不同。而乡职也有「地方总甲、保正、乡约、甲长、墟长、社长等，一切在官人役」70等称谓。乡职是由有地方遴选保举产生，所以「乡约、地方等役，类由本乡本里之民，保送佥充」71，其职责是「责令乡约、甲长等，教化番民，稽查奸匪」。72

\BookSubheading{（一）乡、保的区划}

明郑时，郑经时以陈永华辅政，将东宁分为四坊，分别是东安、西定、宁南、镇北坊，坊置签首，理民事。东宁以外分为三十四里，置总理等乡职，73这是台湾有乡保区划的开始。

清领后在台湾县下区划为十五里、四坊，凤山县下区划为十里、三保、五庄、

\vfill
\begin{BookNotes}
\BookNote{63.}{同上注，页572-573。}
\BookNote{64.}{同上注，页595。}
\BookNote{65.}{同上注，页588。}
\BookNote{66.}{《清奏疏选汇》（南投：台湾省文献委员会，1997年），页77。}
\BookNote{67.}{《清高宗实录选辑》(南投:台湾省文献委员会，1997年），页100。}
\BookNote{68.}{丁曰健，《治台必告录》（南投：台湾省文献委员会，1997年），页349。}
\BookNote{69.}{萧一山，《清代通史·第四篇》(台北:商务印书馆，1976年)，页639。}
\BookNote{70.}{昆岗等着，《清会典事例》（北京：中华书局，1991年），页262页。}
\BookNote{71.}{清高宗敕撰，《清朝文献通考》（台北：台湾商务印书馆，1987年)，21卷，页5045。}
\BookNote{72.}{清高宗敕撰，《清朝文献通考》（台北:台湾商务印书馆，1987年），190卷，页5030。}
\BookNote{73.}{连横，《台湾通史》（台北:众文出版社，1979年)，页555-556。}
\end{BookNotes}

% 原书第 452 页（PDF 第 465 页）
\BookPage{452}{465}
\noindent 十二社，诸罗县下区划成四里、十四庄、三十三社。

台湾县四坊是东安坊、西定坊、宁南坊、镇北坊。十五里为武定里、永康里、广储东里、广储西里、长兴里、新丰里、归仁南里、归仁北里、永丰里、保大东里、保大西里、仁德里、仁和里、文贤里、崇德里。

凤山县十里为新昌里、永宁里、依仁里、长治里、维新里、仁寿里、嘉祥里、观音山里、淡水港东里、淡水港西里。三保为土墼埕保、安平镇保、喜树仔保。五庄为半屏山庄、兴隆庄、赤山庄、大小竹桥庄、凤山庄。十二社为上淡水社、阿猴社、搭楼社、大泽机社、下淡水社、力力社、茄藤社、小琉球社、放索社、加六堂社、郎娇社、卑南觅社。

诸罗县四里是新化里、善化里、安定里、开化里。十四庄是新旧咯庄、大奎璧庄、下茄苳庄、井水港庄、鹿仔草庄、龟佛山庄、南势竹庄、大圻田庄、龟仔港庄、糠榔庄、诸罗山庄、打猫庄、他里雾庄、半线庄。三十三社是新港社、目加溜湾社、萧笼社、麻豆社、大武垄社、大杰巅社、倒咯社、诸罗山社、阿里山社、打猫社、他里雾社、猴闷社、柴里斗六社、奇冷岸社、大居佛社、西螺社、麻芝干社、南社、二林社、大突社、马芝遴社、阿束社、东螺社、大武郡社、半线大肚社、南北投社、麻务栋社、沙辘牛骂社、崩山社、新港仔社、竹堑社、南嵌社、上淡水社、鸡笼社、木武郡赤嘴社、水沙连思麻丹社、麻咄目靠社、挽鳞倒咯社、狎里蝉恋蛮社、干那雾社。74

以上是康熙领台之初台湾、凤山、诸罗三县的乡级区划情形，有坊、里、保、图、庄、社等区分。到雍正时，从诸罗县再分出彰化县及淡水厅，彰化县区划成十保，一百一十庄，淡水厅初设二保，三十五庄。

彰化县有半线保：下辖有登台庄、快管庄、柴坑仔庄、恶马庄、大吼庄、加黎庄、阿东庄、砖仔庄、本县庄、南北投庄、柳树滥庄。猫雾栋保：下辖有义学庄、王田庄、龟山庄、加投庄、水里庄、沙辘庄、牛骂庄、横山子庄、岸里社口庄、枫树脚庄、北势庄、蓝张兴庄、凉伞树脚庄、新兴庄、乌日庄、镇平庄、刘厝庄、阿里史

\vfill
\begin{BookNotes}
\BookNote{74.}{高拱干，《台湾府志》（台北:台湾银行经济研究室，1960年），页35-37。}
\end{BookNotes}

% 原书第 453 页（PDF 第 466 页）
\BookPage{453}{466}
\noindent 庄。燕雾保：下辖有垓庄、秀水庄、佳锡庄、买林子庄、镇平庄、内湾庄、白沙坑庄。马芝遴保：下辖有马明山庄、埔盐仑庄、埔盐庄、大滥庄、南势庄、北势庄、大有庄、打廉庄。东螺保：下辖有旧社庄、三条圳庄、打马辰庄、侯心霸庄、大段庄、十张犁庄、兴化庄、睦宣庄、浦心庄、眉里庄、埤头庄、斗六甲庄、麻园寮庄大武郡保：下辖有旧社庄、仑仔庄、崁顶庄、枋桥头庄、陈厝庄、红毛社庄、油车店庄、火烧庄、滥港东庄、西成庄、柴头井庄、马光厝庄、新庄子庄、卓乃潭庄、桥头庄。西螺保：下辖有样子脚庄、大加冬庄、鹿场庄、大北园庄、何厝庄、埔心庄、下滥庄、张厝庄、加冬仔庄。布屿禀保:下辖有二仑庄、港尾庄、沟子墘庄、南势底庄、马攻厝庄、龙岩厝庄、埔姜仑庄、大坪庄、大庄、萧厝庄、猫儿于庄。二林保：下辖有大突庄、旧社庄、火烧庄、七夕厝庄、火烧厝庄、西势庄、垓庄、过沟子庄、文厝厝庄、西庄、芦竹塘庄、土库庄。深坑子保：下辖有深坑子沟庄、大排沙庄、丈八堵庄、土垄厝庄、犁头厝庄、八人庄。

淡水厅有淡水保：下辖有八里坌庄、沪尾庄、大屯庄、竿蓁林庄、关渡庄、北投庄、八芝连林庄、奇里岸庄、瓦笠庄、兴仔武湾庄、大佳腊庄、圭母子庄、大湾庄、水兴庄、兴直庄、加里珍庄、摆接庄、山脚庄、八里坌仔庄、海山庄、坑子庄、虎茅庄、奶笏庄、涧仔历庄、甘棠庄。竹堑保：下辖有南庄、北庄、芝巴里庄、大溪墘庄、翠丰庄、猫儿碇庄、中港庄、永安庄、猫盂印斗庄、吞霄庄。75

\BookSubheading{（二）保甲编立}

清初「户给印信纸牌一张，书写姓名、丁男口数于上，出则注明所往，入则稽其所来」。76并以「十户为牌（奇零散处通融编列），立牌长，十牌为甲，立甲长，十甲为保，立保长，限年更代，以均劳逸」。77康熙年间台湾的保甲编立情形是：

台湾县：东安坊一百三十甲、西定坊八十八甲、宁南坊四十四甲、镇北坊七十五甲、武定里四十六甲、永康里上保二十甲、永康里下保十八甲、长兴里上保十三甲、

\vfill
\begin{BookNotes}
\BookNote{75.}{刘良璧，《重修台湾府志》（台北:台湾银行经济研究室，1977年），页67-69。}
\BookNote{76.}{高宗敕撰，《清朝文献通考》（浙江古籍出版社，1988年），页5051。}
\BookNote{77.}{徐栋，〈保甲书〉，《古代乡约及乡治法律文献十种》（黑龙江人民出版社2005年），页20。}
\end{BookNotes}

% 原书第 454 页（PDF 第 467 页）
\BookPage{454}{467}
\noindent 长兴里下保十二甲、归仁里南保十八甲、归仁里北保十七甲、仁和里一图十四甲、仁和里二图十一甲、文贤里一图九甲、文贤里二图八甲、崇德里一十八甲、新丰里二十四甲、保大东里二十甲、保大西里四十甲、广储东里一十六甲、广储西里一十九甲、大目降庄二十八甲、仁德里南保一十六甲、仁德里北保一十八甲、永丰里一十七甲。

凤山县:土墼埕保二十甲、安平镇保二十四甲、喜树仔保三甲、新昌里七甲、永宁里八甲、依仁里一图五甲、依仁里二图一十甲、长治里一图一十二甲、长治里二图一十六甲、维新里九甲、仁寿里二十一甲、嘉祥里一十八甲、观音山里二十一甲、半屏山庄一十二甲、兴隆庄二十三甲、赤山庄九甲、大小竹桥庄二十四甲、凤山庄二十四甲、淡水港东里一十八甲、淡水港西里一十八甲。

诸罗县：新化里东保二十甲、新化里西保一十九甲、善化里东保一十八甲、善化里西保一十九甲、开化里赤山保二十五甲、开化里茅港尾保一十五甲、开化里佳里兴保一十六甲、安定里东保二十四甲、安定里西保二十三甲、新旧咯庄一十三甲、大奎璧庄一十一甲、下茄苳庄一十一甲、井水港庄一甲、鹿仔草庄二甲、龟佛山庄二甲、南势竹庄一甲、大坵田庄二甲、龟仔港庄一甲、糠榔庄一甲、诸罗山庄一十甲、打猫庄二甲、他里雾庄一甲、半线庄二甲。78

综上所述康雍干时期台湾的乡治组织有坊里保，再次有图、庄，下有牌、甲，这三层组织在康熙时就已建构完成。唯坊仅在台湾府城所在的台湾县有之，图也仅在凤山现出，里则仅限台湾、凤山二县，保、庄、甲则为普遍存在。另甲的编成，可在坊、里、保、图、庄之下均可，其规模大小，有非常大的弹性，一般是数十甲，但也有一百多甲，例如东安坊就高达一百三十甲，也有一庄仅有一甲。如井水港庄、南势竹庄、龟仔港庄、糠榔庄、他里雾庄均仅有一甲。

\BookSubheading{(三）康雍干时期乡治的运作}

康雍干时期在台湾西部平原及浅山地区设有四县一厅，并划分坊里保，再次有

\vfill
\begin{BookNotes}
\BookNote{78.}{刘良璧，《重修台湾府志》（台北:台湾银行经济研究室，1977年)，页44-45。}
\end{BookNotes}

% 原书第 455 页（PDF 第 468 页）
\BookPage{455}{468}
\noindent 图、庄，下有牌、甲，三层乡治组织，而因「台郡各厅、县向有里差、社差。里差则派定里分，名为催粮，无论开征与否，整年住居该里，与乡保管事探买纳粮，遇事科派。社差则分管社分，名为催纳各番丁饷，亦无论开征与否，整年住居该社」。79又其运作情形，现略述于后。

1、康熙时乡治

在康熙时内地有计户起丁之例或计田起丁之例。然必分别里分，由里长、总甲统理；五年一推收，十年一编审。台湾田园则归之管事、人丁归之保长，就里之大小，或一人、或二三人，终身不改其役。80

初时台湾未有城郭，兵民杂处。为管理的方便。以十家为一甲，设有甲长。十甲为一保，立有保长。其十家、十甲内，有户婚、田土、细事兴讼到官，即差本甲长、本保长勾摄对质，省差人下乡之扰。有来历不明及非为发觉者，甲长、保长不先举报，讯明一体治罪。其甲长、保长须十家、十甲内公举殷实老成之人充当，断不许无赖棍徒包揽生事。81也因台湾「一里一对差，官票在手，无肯空手徒走之理。一里一管事，派歛佃户，未必无指一科十之弊」。82

在番社则设土官「土官有正副，大社五、六人，小社三、四人，各分公廨，有事则集众以议」。83乡职也要应付地方官员的需索，例如豹皮及斑点小鹿皮，文采可观，且无市。每地方官寻觅，必骚动通事、土目，是以有「地方产一异物，便多一累」的说法。84

乡职要负责来往官差食宿，由郡治北至鸡笼旅店、公馆不多，「文武衙门差遣往来，悉主于保长、通事、塘汛，入门则酒食相饷；开台以来相沿为例，不必需索也。

\vfill
\begin{BookNotes}
\BookNote{79.}{《福建省例》（台北:台湾银行经济研究室，1964年)，页849-850。}
\BookNote{80.}{陈文达，《台湾县志》（台北：台湾银行经济研究室，1961年），页229。}
\BookNote{81.}{陈殡，《陈清端公文选》（台北:台湾银行经济研究室，1961年)，页10。}
\BookNote{82.}{同上注，页7。}
\BookNote{83.}{黄叔璥，《台海使槎录》（台北:台湾银行经济研究室，1957年)，页148。}
\BookNote{84.}{周钟瑄，《诸罗县志》（台北:台湾银行经济研究室，1958年)，页292。}
\end{BookNotes}

% 原书第 456 页（PDF 第 469 页）
\BookPage{456}{469}
\noindent 保长、通事之供应既烦，自不得不有所科敛」。85

2、雍正时乡治

雍正四年（1726）斗六东埔地方民人被生番所杀，拾有番弓一把、箭五枝、番牌一面，令通事、土官等看认。为防生番再度杀人并令兵役、土官、乡保等加谨巡防守御。也传同土目、通事、乡保、业户立表定界，约束民人不得越界。86

雍正五年（1727）浙闽总督高其倬奏说:「番人焚杀一节，是因番社俱有通事，通事刻剥，番人愤怨之极，遂肆杀害。已通令各县严行查处通事，不许刻剥番人，胥役不许需索通事。臣亦详行细访，若各社可以不用通事，臣即尽行革除，更为清楚」。87

「编立保甲止是大概」，88所以雍正十一年（1733）福建总督郝玉麟条奏:台属田地、丁口渐增，请将各处庄民编造保甲。其各番社乡保暨通事，俱严加选择，勿令奸匪充当。89并议准：台湾多流寓之人，除立有产业兼有父母、妻子者，自有十家甲牌可稽。其并无家室产业如佃户、佣工、贸易之类，取具业主、房主、邻佑保结，附于甲牌之末汇报，以备稽察。90

3、乾隆时乡治

乾隆四年（1739）巡视台湾御史杨二酉奏台地兵民安帖情形说「保甲之法，行于台地更宜。现在切谕地方员弁行之，颇有成效」。91对新辟地方「耕种男妇编立保甲，设隘防守，不时稽察」，92对内地来台就养者，「准其给照来台，入甲安

\vfill
\begin{BookNotes}
\BookNote{85.}{同上注，页292。}
\BookNote{86.}{郑喜夫选辑，《雍正殊批奏折选辑》（台北：台湾银行经济研究室，1972年)，页35。}
\BookNote{87.}{同上注，页141-142。}
\BookNote{88.}{同上注，页113。}
\BookNote{89.}{《清世宗实录选辑》（南投:台湾省文献委员会，1997年)，页44。}
\BookNote{90.}{《清会典台湾事例》（台北:台湾银行经济研究室，1966年)，页146。}
\BookNote{91.}{《清高宗实录选辑》（南投:台湾省文献委员会，1997年)，页17-18。}
\BookNote{92.}{同上注，页111。}
\end{BookNotes}

% 原书第 457 页（PDF 第 470 页）
\BookPage{457}{470}
\noindent 插。93

乾隆十年（1745)福建布政使司布政使高山上建议:援照川、广苗疆土司之例，令该地方官于众社土目中择其老成诚实、才具明干者数人，呈报督、抚及巡台御史秉公验看，会奏请旨，颁土司职衔，量与顶戴；令其分管各社内番众，统辖生番。凡有番黎争夺戕杀等事，即移令土司分别惩治协拏。设立之后，果能使番众相安，三年无过，该地方官详报具题，量加奖赏，以示鼓励。倘有勾通汉奸生事扰民者革职，另行选举。94

乾隆三十年（1765)奏准：台湾一郡，令地方官编排保甲，即于现在居民内查有人材出众、实系随征出格效力者，移取印、甘各结，给予外委功加札付，委为社长；阖社民人令其约束，统属地方文武管辖。三年之内果能约束有方、民番和辑，该管官出具考语，保送拔用。95

乾隆三十四年（1769）为实力严查、申明禁限，通饬各属将福建、广东赴台人民严行禁止，毋许私行逗遛一人。其现今流寓在台者，虽不能尽行驱回内地；而编设保甲，互相觉察，勿任匪类得以藏奸。96

乾隆五十三年（1788）「只身远渡与挈眷同来之内地民人，应由地方官查明给照、移咨台湾入籍，按户编甲。」97

在《岸里大社文书》也可以见到乾隆二十九年（1764）猫雾抹收取大租谷数清册（如图1），98具名的有通事及社主土官，另在乾隆二十九年（1764）的《岸里大社文书》有乡长、保长、通事、土目、甲头、乡保甲、乡保等乡职名称出现（如图2）。99

乡职有查报不法、凶案或无名尸等职责并呈请报验，如乾隆三年（1738)八月

\vfill
\begin{BookNotes}
\BookNote{93.}{《清世宗实录选辑》（南投：台湾省文献委员会，1997年)，页49。}
\BookNote{94.}{《清奏疏选汇》（南投:台湾省文献委员会，1997年)，页41。}
\BookNote{95.}{《清会典台湾事例》（台北:台湾银行经济研究室，1966年)，页31。}
\BookNote{96.}{《清世宗实录选辑》（南投：台湾省文献委员会，1997年），页199。}
\BookNote{97.}{《清高宗实录选辑》（南投:台湾省文献委员会，1997年)，页621。}
\BookNote{98.}{岸里大社文书，AL00973-000。}
\BookNote{99.}{岸里大社文书，AL00958-128-01。}
\end{BookNotes}

% 原书第 458 页（PDF 第 471 页）
\BookPage{458}{471}
\BookArt{p471-01.jpg}{0.990}{0.306}

\noindent 二十四日，「据竹堑保长陈朝锦禀报:咸水港下海边有一男尸，尸首已经朽烂，理合报明乞验」。100又乾隆七年（1742）「据保长陈朝锦、店邻王恭禀报:初七日，保内王珠与黄宽戏谑黄宽将王珠撞伤等情到职」，101后王珠医治不痊身死。乾隆十二年（1747）有民人唐廷章被抢，乡职猫雾栋西保乡保长郑永瑞、姚文远、管事陈国用立即将前情禀同到县。102乾隆十八年（1753）彰化县民赵悻树旗陷害赵新一案，起因赵悻屡被欺凌，密商同居王兴树旗泄忿，经乡保禀报查缉。103乾隆五十一年（1786)嘉义地区耕牛被偷，嫌犯被捉，等候保长回来解究。104乡职有一定的司法权，对窃取耕牛大案，有解究权限。乡职还要负责社会救济「流丐夜宿庙观，流落疲病，卒于路旁。市民共施棺木、地保为其收埋」o 105

\BookSubheading{（四）清代台湾乡治领导群体}

清代台湾乡治领导群体，本文将没有科举功名的铺户、垦首也包含在领导群体内，这是考量清代台湾是拓垦社会，文教发展较慢，科举并不发达，反而是拓垦豪强

\vfill
\begin{BookNotes}
\BookNote{100.}{《台案汇录己集》（南投:台湾省文献委员会，1997年），页149。}
\BookNote{101.}{同上注，页185-186。}
\BookNote{102.}{同上注，页194。}
\BookNote{103.}{《清高宗实录选辑》(南投:台湾省文献委员会，1997年)，页101。}
\BookNote{104.}{《台案汇录己集》（南投:台湾省文献委员会，1997年），页339。}
\BookNote{105.}{黄典权，《台湾南部碑文集成》（南投:台湾省文献委员会，1994年)，页398。}
\end{BookNotes}

% 原图说　图1　岸里大社文书
% 原图说　图2　岸里大社文书

% 原书第 459 页（PDF 第 472 页）
\BookPage{459}{472}
\noindent 在地方有强大影响力，若单以士绅为唯一乡治领导群体，在清代台湾社会是有严重失真的事实。现将清代台湾乡治的领导群体，依性质分成六个群体加以论述：

1、士绅群体：在职、候选或退休的大小官吏、进士、举人、贡生、例贡、生员

\begin{BookQuote}
以及监生等。当时，台湾士绅不但能和官吏直接往来、出席地方重要庆典、豁免徭役、参与地方公务、履行独特诉讼程式等法律特权，并在社会享有较高的声望，同时也拥有较多的财富。
\end{BookQuote}

2、乡职群体：乡庄组织中的总理、董事、乡保长、街庄长、联甲头人、签首、

\begin{BookQuote}
此类人士为乡庄组织的任职人员，可视为乡里基层组织上的核心领导人物，其地位来自地方推举而得到官方承认。而总理、董事经常是在拓垦早期垦民行自推举出来办理地方公务的自治领袖。
\end{BookQuote}

3、垦首群体：垦首、结首、佃首、隘首和管事、埤圳长、埤圳总理等。垦首是

\begin{BookQuote}
开垦活动的领导者，当然称为地方的领袖，由其出面与官方或民番接洽，这类人是多半没有科举功名。垦首负责规划垦务、筹集资本和保卫垦民安全，官方也给予承认授权，责其管辖垦界之内的行政事务，维持地方秩序。结首是拓垦制主要的领导者，尤其具结于官，约束众佃。开垦进行时，官方或民间为防止番害，就于与番界处设立垦隘，推举隘首，负责防御生番，并负责其地保安事务，有时隘首兼为垦首。管事为垦首的经理人，以管理垦务租务、埤圳长是水利组织的负责人，埤圳总理，是水利组织的经
\end{BookQuote}

理人。

4、铺户群体：郊商、铺户：清代台湾由于商业发达，一些重要的都市都有行郊

\begin{BookQuote}
出现。所谓郊即指行商为互助合作维护成员利益的组织，共分两类：一是以同业商人组成，如糖郊；另一种则是由同一贸易地商人组成，如泉郊等。郊商由于控制台湾的进出口贸易，转而对社会有很大的影响力，逐渐崛起领导公务。
\end{BookQuote}

5、耆老群体：乡宾、耆老、名士硕儒等，此类人士在社会中领导地位的建立，


% 原书第 460 页（PDF 第 473 页）
\BookPage{460}{473}
\begin{BookQuote}
是其有突出的德行与学识，故其社会领导地位的取得在名不在利。乡宾有两类，一是品德高尚的士绅，另一为年高德劭者。康熙末年，台湾即有乡宾的记载，澎湖亦有乡宾出面领导地方的情形；而耆老只要年高德邵为众人所推戴即可。名士硕儒即指社会上因德行高超、学识出众而声望较高、影响力较大的儒士。
\end{BookQuote}

6、其他群体：其他领导群体大致分为三个系统：一是临时公务系统，二是宗族

\begin{BookQuote}
系统，三是寺庙系统。
\end{BookQuote}

\BookSubheading{(1）临时公务系统：}

\noindent 筑城总理、董事、义首：此类人是多是因某项特殊临时姓性事务需要董其事者，经由地方人士推举或政府指派产生，和乡庄组织中的任职人员性质不同。总理、董事往往由地位较高、势力较大的人充任，为地方重要领导人物之一。义首，或称义民首，一般是指出面帮助政府平乱的地方有力人士。

\BookSubheading{（2）宗族系统（族长或族正）:}

\noindent 台湾在清代，政府亦曾下令大族或大姓推举族长或族正，但并不普遍。族长或族正大抵为大姓中辈份较高，又年高德劭之人，故具有极大的影响力。

\BookSubheading{（3）寺庙系统（首事、总理，董事）：}

\noindent 戴炎辉将台湾的寺庙分为官庙与民庙两大系统，民庙又可分为私庙与公庙，公庙又可分为同业公庙、同籍公庙与地方公庙106。而以地方公庙数量最多也最重要，而大多数学者的研究也集中在地方公庙。台湾寺庙是乡治头人拓展势力，支配地方的管道，而诸如做为地方官吏咨询忠告角色、办理地方公共事务与社会福利、领导地方文教、领导团练自卫武力或助官平乱等，经常是经由寺庙来运作。107所以寺庙的领导人首事、总理，董事与

\vfill
\begin{BookNotes}
\BookNote{106.}{戴炎辉，《清代台湾之乡治》(台北:联经出版事业公司，1979年初版)，页38。}
\BookNote{107.}{卓克华，《从寺庙发现历史一台湾寺庙文献之解读与意涵》（台北：扬智文化事业公司，2003年），页8。}
\end{BookNotes}

% 原书第 461 页（PDF 第 474 页）
\BookPage{461}{474}
\begin{BookQuote}
地方领导人大多重叠，对地方公共事务当然有举足轻重的影响力。
\end{BookQuote}

\BookSubheading{(五)清代台湾乡治的头人}

台湾将地方的领导人称「头人」，头人不是专有名词，是混称，总理、街庄正、仕绅耆老都被称为地方头人。108因台湾五方错处，事杂言庞，是以有头人之设。向例惟家道殷实、素行端谨者，方准举充。官课以考成并赏罚之，故于地方有益。但头人也以官戳为护符，以文檄为奇货，竭良善之脂膏，适以饱豪强之囊橐。即有贤明有司，也多坠其术中而不悟。109

头人官府有给谕戳，所以才会「以官戳为护符，以文檄为奇货」，而查考淡新档案，官府给戳记的是总理、董事、街庄正、地保、族长等。所以戴炎辉认为头人有广狭二意，广意指一切领导人，包括街庄组织负责人、绅衿、耆老等，狭义则则指总理、董事、街庄正、地保、族长等。110

1、狭义的头人：有职衔戳记，也有任命公文，公事上有特定职权者。

\BookSubheading{（1）总理（签首）：}

\begin{BookQuote}
总理是清代台湾最重要的乡职「为一堡之正」、111「为街庄之首」，112是「地方自治团体的首席，且为其执行人」。113官府对总理有「给戳给谕之举」，114而遭到退革的总理要将「前领谕戳禀缴察销」。115总理的居所被称为公馆，116门前还挂有两个大灯，灯上也写有总理两个大字，117总理俨
\end{BookQuote}

\vfill
\begin{BookNotes}
\BookNote{108.}{戴炎辉，《清代台湾之乡治》（台北:联经出版事业公司，1979年初版)，页178。}
\BookNote{109.}{吴子光，《台湾纪事》(台北:台湾银行经济研究室，1959年)，页67。}
\BookNote{110.}{戴炎辉，《清代台湾之乡治》（台北，联经出版事业公司，1979年初版），页38。}
\BookNote{111.}{「淡新档案，第12229-5号。}
\BookNote{112.}{「淡新档案」，第12203-11号。}
\BookNote{113.}{戴炎辉，《清代台湾之乡治》（台北：联经出版事业公司，1979年初版），页21-22。}
\BookNote{114.}{「淡新档案」，第12301-5号。}
\BookNote{115.}{「淡新档案」，第12238-3号。}
\BookNote{116.}{白长川，《罗东历史地名寻根》（宜兰：罗东镇公所，2003年），页17。}
\BookNote{117.}{王凯泰，〈台湾杂咏续咏〉《台湾杂咏合刻》（台北：台湾银行经济研究室，1959年），页42。}
\end{BookNotes}

% 原书第 462 页（PDF 第 475 页）
\BookPage{462}{475}
\noindent 然是一位官员，他的住处俨然是地方的官署。签首与总理相当，明郑时「分都中为四坊」，「坊置签首，理民事，制鄙为三十四里，置总理」118。明郑在东都城内设签首，城外设总理，以处理地方事务，但清代签首制只实行在台南一带，在北部的淡水厅并没有发现签首制，而在台南一带也是签首和总理并行，在「严禁窃砍竹城碑」记有「台属永康里龙山寺境总签首王文忠、总理汪博厚」119等，请求台湾兵备道吴大廷，示禁窃砍东门外竹城的竹围。《安平县杂记》在「捕缉官名」列有:「安平县、捕厅、巡检、保甲委员、文口委员、皂役、快役、局勇、地保、总理、乡长、总签首、签首、壮丁、更夫、庄耆、中营中军、城守、左军、右军、水师、营勇、汛官、汛兵、武口汛弁·」120。地保、总理、乡长、总签首、签首等并列为补官，负责地方事务。

\BookSubheading{(2)董事：}

\noindent 淡水同知娄云在「庄规四则」第一则就说：「各庄向设总理、董事、庄正、庄副，官给札谕戳记约束庄众」，121清代台湾街庄都设有董事，董事作为乡职，是在协助总理，所以每一街庄经常是总理一人，董事数人。吞霄街总理梁壬生，在禀举陈存仁担任董事的文中说「缘生等该处保内，原有董事协办公务，前蒙示革，为蒙复充，恐地方有事，助治无人」，所以禀举陈存仁，以「帮同总保、街庄正副约束庄规」。122董事在地方是助治、帮同的辅佐性质，不是执行者，没有单独行为权力。

\BookSubheading{(3)街庄正：}

\noindent 清代台湾各街庄，设有街庄正副，由总理与地方绅耆推举，经认充后由官府发给戳记。其主要职务是管理保甲的编成，调处街庄内的纠纷。

\vfill
\begin{BookNotes}
\BookNote{118.}{连横，《台湾通史》（台北：台湾银行经济研究室，1961年)，页555。}
\BookNote{119.}{黄典权辑《台湾南部碑文集成》（台北：台湾银行经济研究室，1968年)，页494。}
\BookNote{120.}{不着撰人，《安平县杂记》（台北:台湾银行经济研究室，1959年)，页44-45。}
\BookNote{121.}{《淡水厅志》（台北:台湾银行经济研究室，1959年)，页388。}
\BookNote{122.}{「淡新档案」，第12203-28号。}
\end{BookNotes}

% 原书第 463 页（PDF 第 476 页）
\BookPage{463}{476}
\BookSubheading{（4）保（乡）长：}

\noindent 乡保长是官府差役的在地化，官府在各乡保派有常驻的差役乡保长，以做为官府耳目，而保长与管事都没有任期，所以说:「今台湾田园归之管事、人丁归之保长，就里之大小，或一人、或二三人，终身不改其役」。123乡保长除在保内人丁查察以外，另一主要工作是负责政令的传导，张贴官府告示是乡保长最主要工作，官府告示交给乡保长后，乡保长要立刻以「满浆实贴」，124张贴在保内街庄。

\BookSubheading{（5）义首：}

\noindent 「清代治理下的台湾是变乱频仍的地区，『三年一小反，五年一大反』殆尽人皆知之俗谚」，125台湾变乱多，官府自朱一贵民变开始就大量运用各地头人为义首，组织义民协助官府压制变乱，所以义首在清代台湾社会是地方有企图人士，想取得头人头衔开了一扇方便的门，所以清代台湾由义首率领的民兵部队有时非常庞大。雾峰林家的林文察由义首转战各地并升到总兵，可以说是「躬膺显职」126，吴沙也是以义首名义号召闽粤移民进入宜兰拓垦。

\BookSubheading{(6)隘首：}

\noindent 清代台湾是属武力拓垦，所以要在拓垦的外缘设隘，隘首的职务是督导隘丁，防守生番，而隘有官隘，有民隘。官隘则额定隘丁若干名，需工食银若干圆，皆由隘首向官支领，而后给与隘丁，无敢透遛者。民隘则取粮于农，以佃田之多寡为等差。凡应隘丁募，先议定工价；倘遇不测，矢无悔意，不过给以殡埋薄赀而已。又隘之险夷无定，有昔险而今夷者，皆当因时制宜，总以防堵备御为上策。故隘首与垦户事不相谋，而弊恒相因。清代台湾的隘务很复杂也很重要，所以隘首的任命非常谨慎，要由「承耕

\vfill
\begin{BookNotes}
\BookNote{123.}{陈文达，《台湾县志》（台北:台湾银行经济研究室，1958年)，页229。}
\BookNote{124.}{「淡新档案，第12507-10号。}
\BookNote{125.}{刘妮玲，《清代台湾民变研究》（台北：国立台湾师范大学历史研究所，1982年），页1。}
\BookNote{126.}{丁日健，《治台必告录》（台北:台湾银行经济研究室，1959年)，页582。}
\end{BookNotes}

% 原书第 464 页（PDF 第 477 页）
\BookPage{464}{477}
\noindent 课地各佃及往山樵探诸人选举」，127并要互相订约以明责任，最后还要经官府任充，给予戳记才能取得隘首资格。而隘首处理隘务以外，也要参与地方事务「护卫庄民」，128公同订定庄规，129推举地方总理。130

\BookSubheading{（7）族长（族正）：}

\noindent 《福建通志台湾府》记有「择其乡之齿长而端悫者立族正及副二人，如古三老啬夫，凡乡有讼事，族人以告族正，小事族正判其曲直，大事则族正自诣县告」，131族正是家族内是非的仲裁者，对怙恶不悛的族人，族长可公开责惩，严重的可逐出本支或联名送官究处，不准他回社。如改过自新，则可以推荐到衙门充当壮勇。132而族正要「编查保甲」，133也要对匪类有举报的责任，若徇私隐满，会受到连坐处分134。

\BookSubheading{（8）土官、土目：}

\noindent 台湾土官初设于荷兰时代，明郑也沿袭荷兰以土官进行部落自治，清朝也沿用这一制度，用土官管理各番社，到了乾隆年间渐以土目取代土官的称呼。土目是由各番社屯目、甲首、耆老、社众推举，经官府验充后发给戳记。土目对内管理社务，约束社众，管理公租，发给口粮。对外代表番社，处理汉番事务135，但不介入汉人街庄有关的公共事务。

\BookSubheading{(9) 通事：}

\noindent 通事设立在明郑时代，清代沿用，「通事，则由社总及土官、土目暨各社、各庄耆老番民选举堪专其任之人，禀官批准充当」136，而当通事「必

\vfill
\begin{BookNotes}
\BookNote{127.}{丁绍仪，《东瀛识略》（台北：台湾银行经济研究室，1957年），页48。}
\BookNote{128.}{「淡新档案」，第12306-1号。}
\BookNote{129.}{「淡新档案」，第12203-7号。}
\BookNote{130.}{「淡新档案」，第12212-5号。}
\BookNote{131.}{《福建通志台湾府》（台北:台湾银行经济研究室，1960年)，页737。}
\BookNote{132.}{徐宗干，〈斯未信斋文集〉《治台必告录》（台北:台湾银行经济研究室，1959年），页366。}
\BookNote{133.}{《清宣宗实录选辑》（台北:台湾银行经济研究室，1964年)，页47。}
\BookNote{134.}{屠继善，《恒春县志》（台北：台湾银行经济研究室，1960年），页132。}
\BookNote{135.}{洪丽完，《台湾古文书》（台中：台中县立文化中心，1996年），页8。}
\BookNote{136.}{不着辑人，《安平县杂记》（台北：台湾银行经济研究室，1959年），页68。}
\end{BookNotes}

% 原书第 465 页（PDF 第 478 页）
\BookPage{465}{478}
\noindent 能知番语及熟悉番社情形，并知详官场礼节者，始可选充其任」。清统治台湾初期以汉人为通事，后渐由土着充任。通事的职责在疏通民番，兼及教化，也办理番饷及赋课差役。而若社群较具规模，社务浩繁，也有增设总通事，以统里社务。137

2、广义的头人：指一切地方领导人，包括狭义的头人。

\BookSubheading{（1） 绅衿：}

\noindent 瞿同祖指出绅衿是「功名持有者」，138生员、监生、贡生、举人、进士都是绅衿也称为绅士。他们在地方受到尊敬与追从，是地方的意见领袖，而绅衿也实际介入街庄事务，充当乡职等职务，乾隆五十一年（1786）任六堆总理的曾中立就是是举人。

\BookSubheading{(2） 耆老：}

\noindent 蓝鼎元认漳泉与客庄居民都是台湾百姓，宜相爱相亲，「或有言语争竞，则投明乡保耆老，据理劝息，庶几兴仁兴让之风。敢有攘夺斗殴，负嵎肆横，本镇执法创征，决不一毫假借。其或操戈动众相攻杀者，以谋逆论罪，乡保耆老管事人等一并从重究处」，13耆老是地方管事，官府对耆老十分依重，但惩处也十分严厉。而耆老是地方的头人，官员到任会先拜访耆老，征询耆老意见，这是一种尊重，也表示耆老在地方的重要性，姚莹也说:「奉檄权判其地，访诸耆老」。140而耆老要参与各种乡饮酒礼、耕耤等典礼，对寺庙重建，清庄联甲等都要参与。

\BookSubheading{（3）郊行、铺户：}

\noindent 郊行是商业组织，是地方财力的主要提供来源之一，而清代台湾郊行除台南以外、遍布鹿港、盐水港、北港、艋舺、新庄、凤山、嘉义等地。141郊

\vfill
\begin{BookNotes}
\BookNote{137.}{洪丽完，《台湾古文书》（台中：台中县立文化中心，1996年），页7。}
\BookNote{138.}{瞿同祖，《清代地方政府》（北京:法律出版社，2003年），页284。}
\BookNote{139.}{蓝鼎元《东征集》（台北：台湾银行经济研究室，1958年)，页81。}
\BookNote{140.}{姚莹〈噶玛兰原始〉《噶玛兰厅志》（台北：台湾银行经济研究室，1957年)，页371。}
\BookNote{141.}{蔡渊絜，《清代台湾的会领导阶层》（台北:国立台湾师范大学历史研究所硕士论文，1980年)，页77。}
\end{BookNotes}

% 原书第 466 页（PDF 第 479 页）
\BookPage{466}{479}
\noindent 行由控制经济，进而对社会有很大影响力，逐渐成为台湾社会领导阶层的重要成员，介入地方的公共事务，推举地方总理，郊户也出任总理。郊行铺户与衿耆一样受到官府的重视，并被赋予推举地方总理的权力，所以郊行由商权，伸展到地方共共事务，成为变相的行政机构，上要接应官谕，下要处理地方是庶务。142

\BookSubheading{（4）垦首、结首：}

\noindent 垦首是向官府申请许可，而后出资招佃开垦，每年向佃户收租一成，其他钱粮正供、番租，还是佃户要完纳，与垦首无干。143因此不肖绅民与地方勾结，知道何处有田可垦，就以垦首名义递禀承包，然后分给垦户自行开垦，「垦首但呈一禀，不费一钱，成熟后，垦首每岁抽租一成，就是所谓的大租」。144官府准许垦首开垦会先给垦谕，等垦透报升再给垦照，垦首领到垦照就表示拥有土地所有权，145可得庞大利益，累积可观财富。

\BookSubheading{(5) 管事：}

\noindent 管事是垦户（业户）派驻在庄的经理，也就是租馆的管事，为垦户管理田产与收取租谷，藉垦户势力，管事在地方有相当影响力，也介入地方事务。官府就常将赋税、差役为由管事办理，「乡推一人理赋税差役，官就而责成之」。146而庄民散处，佃户居多，管事穿梭期间，「一庄之中，孰奸孰良，管事多能辨别」。147因管事是业主的雇佣负责管庄，代收租息，乡间来往，闻见最真，是官府重要的耳目。

\BookSubheading{（六）乡职的职权}

\noindent 明郑时规定「凡人民的迁徙、职业、婚嫁、生死，均要报于总理，总理再汇报于

\vfill
\begin{BookNotes}
\BookNote{142.}{卓克华，《清代台湾的商战集团》（台北:台原出版社，1990年初版)，页158。}
\BookNote{143.}{《台湾私法物权编》（台北:台湾银行经济研究室，1963年)，页37。}
\BookNote{144.}{刘铭传，《刘壮肃公奏议》（台北:台湾银行经济研究室，1958年)，页304。}
\BookNote{145.}{《清代台湾大租调查书》（台北:台湾银行经济研究室，1963年），页13。}
\BookNote{146.}{余文仪，《重修凤山县志》（台北：台湾银行经济研究室，1957年)，页98。}
\BookNote{147.}{《台案汇录己集》（台北:台湾银行经济研究室，1964年)，页291。}
\end{BookNotes}

% 原书第 467 页（PDF 第 480 页）
\BookPage{467}{480}
\noindent 官。而总理的工作还有「考其善恶，信其赏罚，劝农工，禁淫赌，计丁庸，严盗贼。而又训之以诗书，申之以礼义，范之以刑法，励之以忠敬」148。乡职总理在地方有非常大的权限。到清代乡职总理，职权更形扩大，也更为明确，现就以一般职权、自治职务、官府委办职务、官给谕帖中规定的职权四方面条例说明。

1、一般职权：

清代台湾乡职一般职权，在《新竹厅志》、《新竹县志》都有论及，近来学者戴炎辉、周宗贤也曾加以讨论，现综合如下：

\BookSubheading{（1）代表街庄维护街庄利益，对外或对官府进行交涉。}

\BookSubheading{(2)调解地方田土户婚雀角等各种纠纷。}

(3)撤回已告官的诉案，或接收官署驳交案件，给予调处。

\BookSubheading{（4）约束庄众，维持地方秩序。}

（5）主动查察地方不法，或被动配合官署查办不法。

\BookSubheading{（6）协助官署缉捕、挐送人犯。}

\BookSubheading{（7）编查保甲，门牌制发管理。}

\BookSubheading{(8）办理团练、冬防，保甲联庄、清庄。}

\BookSubheading{（9）发动地方迎送上官，并予接待。}

\BookSubheading{（10)管理地方公共事业。}

\BookSubheading{(11)地方公务经费的劝募与筹措。}

\noindent 2、自治职务：

\BookSubheading{（1）教化及约束庄民，取缔不法，维持境内安全。}

\BookSubheading{(2）防止生人入庄，补挐来侵匪类。}

(3)接受庄民斗竞、争执等纠纷投诉，并给予排解、仲裁。

（4）负责董事、街庄正、隘首等乡职的禀举与斥革。

（5)负责地方寺庙、道路、桥梁、义塾、义渡、义仓、义学、一冢、修建与管理。

\vfill
\begin{BookNotes}
\BookNote{148.}{连横，《台湾通史》(台北:众文出版社，1979年)，页555-556。}
\end{BookNotes}

% 原书第 468 页（PDF 第 481 页）
\BookPage{468}{481}
\noindent 3、官府委办的职务：

官府委办职务，约可分为三类；一是行政类，二是司法类，三是民事类。官府委办职务是总理最繁重的工作，也是非执行不可的工作，所以非身强力壮，是不能胜任，这也是清代台湾总理等乡职，都由四、五十岁，正当壮年者出任的原因所在。

（1）行政类：官府谕告的传达与收管、赋税的催课、户口保甲的编查、清庄联

\begin{BookQuote}
庄团练、公差派遣与督导、失物的保管、垦隘的督办。
\end{BookQuote}

(2）司法类：民、刑事案件的禀报、查覆、勘验、人犯的追捕解送、证人的传

\begin{BookQuote}
讯、裁判的执行等，地方口角、斗殴、骚扰、强夺、掳掠的谕止，路尸的处理等。
\end{BookQuote}

（3）民事类：总理等乡职是乡庄的重职，各类民事仲裁，或为公亲，或为知

\begin{BookQuote}
见，要在字据、契约、阄书上签署用印。清代规定房地买卖契约要乡约署名才能生效，立契后乡约还要负责催课，不然要受罚，所以清代买卖契约后面盖有「买卖房地总以乡约书名为凭，立契后该乡约即催令赴宪投税。倘迟至十日不税者，或被告发，或经访闻，除业主照漏税例治罪外，将乡约一并究处」149。而这与戴炎辉认为立契，为公亲或知见「非总理的职务，乃附随的服务」150似乎稍有出入。而清代台湾契字有「红契」与「白契」的区别。所谓红契是指契字依官府规定向地方衙门缴税登录有案，这一类契字有已缴契税证明的「契尾」，并盖有红色的关防印信，所以也称为「印契」。「红契」是合法契约保障性较高，但要缴纳一定契税（约为田价的3\%），大多数人为避免这笔花费，在签订契约时，只是聘请双方信任的仲介人士和相关亲属见证签字，就算完成手续，这些地契即是民间所谓的「白契」。白契是违法的，如被告发或访闻到，要被依漏税例治罪，总理也要一并
\end{BookQuote}

究处。

\vfill
\begin{BookNotes}
\BookNote{149.}{田涛、宋格文、郑秦编，《田藏契约文书粹编（一）》（北京：中华书局，2001年)，页28。}
\BookNote{150.}{戴炎辉，《清代台湾之乡治》（台北:联经出版事业公司，1979年初版)，页32。}
\end{BookNotes}

% 原书第 469 页（PDF 第 482 页）
\BookPage{469}{482}
\BookSubheading{(4)谕帖明定的职权：}

\noindent 在总理等乡职验充后，厅、县官会发给谕帖，也就是乡职的人事命令，在谕帖上，厅、县官会对总理等乡职的工作有一些规定和要求，当然也赋予这方面的职权，主要内容有:「如遇地方大小事故，应即出为排解。现值冬令，筹防关要。倘有盗警，尤当督率壮丁护救，并围拏解办」、151「赶紧清庄、联庄，于扼要路径，搭盖更楼，派拨庄民支更把路；实力奉行」、152「如遇地方有口角细故，务须出为排改息事。倘有盗警，尤宜督率壮丁，晓夜巡防」、153「务须勤慎奉公，整顿庄规，约束乡愚；安分守业，不得恃强逞斗、并容留匪类潜入境内。一面督同乡保，办理清庄事宜，严密稽查，巡缉盗匪；如有真赃正贼，务即严拏解案，毋稍纵延。至于庄内如有口舌细故，亦即排难解纷，妥为调处，不得唆弄是非，借端需索」、154「嗣后保内民人，如有田土纠葛，以及口角细故，即行秉公，妥为解散，毋任争斗。倘有命盗重案，随时禀报；一面协差，拏犯𬘡送。至保甲一项，尤关要紧，即速按户，查明人口，塡注清册，呈候塡给门牌悬挂」、155「如遇街庄大小事故，务须出为排解息事，毋许庄民聚赌滋端。倘有盗警，尤宜督率庄丁，协力围拏；并于冲要处所，晓夜梭巡，毋使盗贼逗遛境内。其有命案重情，并当据实禀报，慎勿包庇隐、以及索诈扰累」156、并一再告诫总理等乡职「凡遇应办之事，均须实力奉行，勿得怠惰玩忽」。交办事项「务须认真妥办，毋得始勤终怠，切勿借端需索」，「倘敢故违、妄作，三尺俱在，决无轻纵」。要总理等乡职「凛之，慎之，毋违」。在谕

\vfill
\begin{BookNotes}
\BookNote{151.}{「淡新档案」，第12223-4号。}
\BookNote{152.}{「淡新档案」，第12203-9号。}
\BookNote{153.}{「淡新档案」，第12218-5号。}
\BookNote{154.}{「淡新档案」，第12214-9号。}
\BookNote{155.}{「淡新档案」，第12228-4号。}
\BookNote{156.}{「淡新档案」，第12219-3号。}
\end{BookNotes}

% 原书第 470 页（PDF 第 483 页）
\BookPage{470}{483}
帖中可清楚看到总理等乡职职权非常重，厅、县官对其等寄望很深，要求也很殷切。

\BookSubheading{四、清代台湾乡治的发展与特色}

\BookSubheading{（一）清代台湾乡职有常设、普设化现象：}

总理是清代台湾普设、常设的乡职，因是普设、常设特别显现出，总理在清代台湾的被需要性与重要性，也因总理的被需要与重要，所以必须普设与常设。在同治十一年（1872）淡水同知周式濂向台湾知府呈报「淡属各保总理、董事清册」（如图3）：

\BookArt{p483-01.jpg}{0.990}{0.568}


% 原图说　图3　淡属各保总理、董事清册

% 原书第 471 页（PDF 第 484 页）
\BookPage{471}{484}
从「淡属各保总理、董事清册」中可知淡水厅二十保当时设置总理情形：

首保本城（竹堑城）有：东门总理：林来泽，西门总理：林益修，西门外总理：陈嘉福，南门总理缺，南门外总理：刘贤文，北门总理悬缺，北门外总理：黄应平。

竹南一保有：中港街总理：叶春魁，后庄总理：徐琳盛，头份保总理：谢焕光，三湾庄总理：黄国霖。

竹南二保有：后垅街总理：杜和安，猫里街总理：谢镇基，六大庄总理：徐佳福，双连潭总理：吴轧德，蛤仔市总理：张焕彩，三叉河总理：杨德友、李德龙，奎隆庄总理：彭继生、黄永芳、吴阿三，铜锣湾庄总理:李沧玉，樟树林总理:吴金生，头二三四湖总理：陈振钢，鸭母坑庄总理：彭清禄。

竹南三保有：吞霄街总理:张凤岐、黄有升，白沙墩庄总理：陈明元，宛里街至房里总理:陈金源，日北山庄总理：陈振文。

竹南四保有：大甲街总理：梁克忠、谢玉麟、王朝栋、王昆岗、蔡学礼，内埔庄总理：陈王爵，墩仔脚庄总理：张贻谋，大安港总理：黄片。

竹北一保有：南势庄总理：陈清辉，牛埔庄总理：吴有义，香山街总理：杨保泰，羊簝虎仔山庄总理：郑德宗，海口庄总理：彭荣春，九芎林总理：叶国钦，北埔庄总理：黎德福。

竹北二保有：新埔街总理：刘秉中，咸菜瓮总理：罗光烈、杨英华，大湖口总理：叶呈华，杨梅坜、头重溪总理：邱殿安，波罗汶庄总理：黄锡兰、张清荣，大溪瀚总理：吴德辉、


% 原书第 472 页（PDF 第 485 页）
\BookPage{472}{485}
\noindent 徐国清。

桃涧保有：总理:邱紫来、吴振生、南崁庄总理:李容裕、霄里庄总理:吴维山、三坑仔庄总理：游建邦、龙潭坡总理：古史荣、张维纲、邱庆福，大姑崁总理：林愈贤、张新潭，大坪顶总理：王和，龟仑口总理：陈奠邦，中坜街总理：悬缺。

海山保有：三角涌庄总理：周孙履、陈玉芳、林志道，潭底庄总理：简开瑞、简国安。

大加蚋保有：艋舺街总理：蔡达琪，奎府聚总理：陈王记，溪州底庄总理：杨逢时，锡口街总理：高永辉、王义源，后山埤庄总理：林彩章，南港庄总理：杜冲霄，新南庄总理：潘文桂。

拳山保有：总理：高宽裕，深坑仔庄总理：黄蟾辉、高文印，木栅庄总理：张援萃，视尾庄总理：高星辉，十五份庄总理：林澄清、林维。

摆接保有：总理：沈陈成，枋桥庄总理：林新传、林永顺，溪州庄总理：陈文忠，暗坑三城总理：林松花。

石定保有：总理：陈明智，暖暖街总理：吴盘枝。

八里坌保有：总理：赖瀛州、卓宗陈、古文华、曾文福。

兴直保有：总理：高玉辉、张云龙、林鸣泉。


% 原书第 473 页（PDF 第 486 页）
\BookPage{473}{486}
\noindent 芝兰一保有：总理：赖文渊。

芝兰二保有：总理：李宴林、高良安、陈世清。

芝兰三保有：总理：陈集贤、沈廷玉、简秉干、陈克铭、朱源清。

金包里保有：总理：李有察。

鸡貂保有：总理：何拱辰。

从清册中发现少有保总理悬缺，全厅悬缺的仅有首保（竹堑城内）南门总理与北门总理。而每一保或每一街庄总理可以是一人或数人，最多的芝兰三保和大甲街各有五名总理。

另在「淡属各保总理、董事清册」公文稿中，周式濂批示有「姓名务须核对底卷，毋任桀错，切切」，157显见清代台湾厅县衙门建有总理名册，所以有底卷可以核对。在光绪十六年（1890）总理蔡文光举充一案中，新竹知县批示:「候饬经承注册，发给谕戳办公」，158更说明总理必须建册列管。也说明总理应是常置，受到衙门的列管。而《新竹县制度考》也详列新竹县辖下各街庄总理，印证总理是县的下级建置，街庄普设：

竹北一堡：

\begin{BookQuote}
油车港（即旧港），彭逢春。香山塘，卓清兰。北埔庄，何廷辉。
\end{BookQuote}

\vfill
\begin{BookNotes}
\BookNote{157.}{「淡新档案」，第12213-4号。}
\BookNote{158.}{「淡新档案」，第12243-1号。}
\end{BookNotes}

% 原书第 474 页（PDF 第 487 页）
\BookPage{474}{487}
\begin{BookQuote}
九芎林，范士珍。树杞林，贝登阙。蔴园庄，陈钦云。
\end{BookQuote}

竹北上二堡：

\begin{BookQuote}
崁头厝，黄瑞云。杨梅坜，邱殿安。大湖口，苏如海。大溪墘（溪南），刘兆荣。大溪墘（溪北），钟廷英。许厝港，林清兰。
\end{BookQuote}

竹北下二堡：

\begin{BookQuote}
新埔街，朱逢光。咸菜瓮，谢有时。坪林上下南片，曾国南。
\end{BookQuote}

竹南一堡：

\begin{BookQuote}
头份街，陈光海。中港街，陈如藩。南隘庄，范克昌。159
\end{BookQuote}

从《新竹县制度考》书中可知，总理是县的建置，而在《恒春县志》也记有「兹恒邑各庄，均已设有总理，或一人、或两人，事适繁多，不能周转；且其中公事，有非总理所能办者，自应另设地保一人以副之」160。

台湾总理的设置从甚么时开始？戴炎辉以淡新档案为依据，推测「地方置总理，依保守看法，始自嘉庆年间」161。戴炎辉的推论成为定论，蔡渊絜等都秉持这一看

\vfill
\begin{BookNotes}
\BookNote{159.}{不着撰人，《新竹县制度考》（台北：台湾银行经济研究室，1961年），页12-14。}
\BookNote{160.}{屠继善，《恒春县志》(台北:台湾银行经济研究室，1960年)，页132。}
\BookNote{161.}{戴炎辉，《清代台湾之乡治》（台北:联经出版事业公司，1979年初版)，页20。}
\end{BookNotes}

% 原书第 475 页（PDF 第 488 页）
\BookPage{475}{488}
\noindent 法，而连横在《台湾通史》书中却说总理在明郑时就开始设置，而且清得台后沿用明郑制度，连横是这样说：

\begin{BookQuote}
经立，委政勇卫陈永华，改东都为东宁，分都中为四坊，曰东安、曰西定、曰宁南、曰镇北。坊置签首，理民事。制鄙为三十四里，置总理，里有社；十户为牌，牌有长；十牌为甲，甲有首；十甲为保，保有长；理户籍之事。凡人民之迁徙、职业、婚嫁、生死，均报于总理。仲春之月，总理汇报于官。考其善恶，信其赏罚。劝农工，禁淫赌，计丁庸，严盗贼。而又训之以诗书，申之以礼义，范之以刑法，励之以忠敬，故民皆有勇知方，此则郑氏乡治之效也，清人得台，沿用其制162。
\end{BookQuote}

连横的说法在高屏六堆地区可得到印证，六堆指的是中堆（竹田乡）、左堆（佳冬、新埤乡）、右堆（美浓镇、高树乡）、前堆（麟洛、长治乡）、先锋堆（万峦乡）、后堆（内埔乡），横跨九个乡镇的大客家聚落，所以他的总理称「六堆大总理」。六堆第一任大总理李直三是在康熙六十年（1721）被推举出来，第二任大总理侯心富雍正九年（1731）继任，第三任大总理曾中立是乾隆五十一年（1786）继任。163

六堆自康熙到光绪一百七十多年间，前后相承，历经十任大总理，而在各堆中也有总理，就以右堆为例。右堆首任总理是美浓的开拓者林丰山，他的任期非常长有四十一年，自乾隆元年到乾隆四十一年（1736-1776）、第二任林长炽乾隆四十一年到乾隆五十年（1776-1785)、第三任林向选乾隆五十年到嘉庆二十五年（1785-1820）。164六堆的总理是六堆地方最高领导人，并率领六堆子弟参加各种战事从朱一贵民变到抗日，可说无役不与。而这并非是六堆独有的现象，在清代台湾汉人移垦初始总理就有常设、普设的现象。

\vfill
\begin{BookNotes}
\BookNote{162.}{连横，《台湾通史》（台北:台湾银行经济研究室，1961年)，页555-556。}
\BookNote{163.}{钟壬寿，《六堆客家乡土志》（屏东:长青出版社，1973年)，页175-184。}
\BookNote{164.}{同上注，页185。}
\end{BookNotes}

% 原书第 476 页（PDF 第 489 页）
\BookPage{476}{489}
\BookSubheading{（二）清代台湾乡职有完备的举充程序}

清代台乡职湾总理等的举充有一定的程式，当街庄总理出缺，地方头人以生员或总理为首，在公所或寺庙集会共推一人为总理，再检附保结状和被举人的任充结状，向厅、县禀举，县令认为不适任，就予驳斥，若符合资格，就传讯受验，讯验后发给谕帖与戳记，总理的举充就分为保举与验充两阶段：

1、保举：

在总理的选举中，参与选举的地方头人，要签具保结状，保结总理可以「实心奉公，不敢玩误」。165如有违误，惟保结人是问，选举要附连带保证责任就称为保举。166

\BookSubheading{(1)选举：}

\begin{BookQuote}
在街庄总理因病故、他迁或斥革出缺时，地方头人以生员或总理为首，在公所或寺庙集会妥议，推选出总理人选后，头人等会「公立请帖」167送交被推举的总理以表示尊重及留有纪录。在同时也要将选举结果向官府禀举，在禀举时也要检附保结状及被举充总理的任充结状，在禀中头人要具名用印，并要说明禀举理由；禀文(如图4)
\end{BookQuote}

光绪二年（1876）九芎林等庄绅衿、垦户、乡耆、佃户、铺户等，禀举姜荣椿为总理的禀文内容：

\begin{BookQuote}
具禀治下竹北一保九芎林等庄监生詹国和、生员魏缵唐、垦户徐荣章、业主潘坤光、门生詹鹏材、职员刘维垣、生员魏盛唐、监生林步蟾、监生曾清澜、职员詹廷安、职员魏祥衢、监生刘耀南、乡耆詹达材、徐富业、邓春松、詹五和、童生黄玉璋、陈鸿宾，暨众绅士、佃铺户人等，为总理悬缺，佥举承充恳恩准给谕戳，以专责成，以靖地方事」。切思：朝廷求硕辅，遴选必得忠孝之英；乡里立正人，办公须择明干之士。缘九芎林、五和庄、尾仔、树
\end{BookQuote}

\vfill
\begin{BookNotes}
\BookNote{165.}{「淡新档案」，第12206-2号。}
\BookNote{166.}{「淡新档案」，第12202-18号。}
\BookNote{167.}{「淡新档案」，第12203-14号。}
\end{BookNotes}

% 原书第 477 页（PDF 第 490 页）
\BookPage{477}{490}
\BookArt{p490-01.jpg}{0.945}{0.377}

\noindent 杞林，外至上下员山、下山、六张犁所有一带地方，自上年开垦以来，前经佥举姜秀銮充当总理，公正勤劳，名誉显赫。迨銮年老，经众再举銮姪姜殿魁接充总理，约束地方，最为安静。业蒙王前宪奖赐匾额「义卫乡闾」四字。及魁故后，经众又举伊子姜荣茂接充总理，亦能急公效义，排难解纷。足见继继绳绳，世济其美。况一保大小三十余庄，地广事繁，现总理彭捷科、徐安邦早已示革，即张林超、刘世珍近亦身亡。纵前举有庄正、乡保诸人，未免疏漏不及；倘有公干，谁为统率。生等忝属生长，念切地方，爰集众议，佥举姜荣椿承充总理，将九芎林等庄，归一专管。查其生平，正大光明，志足以绳其伯祖，才足以继其父兄，兼有身家，众皆悦服。唯蒙准给谕戳承充，俾额不多设，人不滥充，庶不负宪恩兴利除弊之至意也。势得连名佥禀，并粘保结。伏乞公祖太老爷好恶同民，恩准姜荣椿承充总理，迅给谕戳，俾一保地方，有专责成，以办公务。斯删除简，而选举当，僻疆幸甚，百姓爱甚，沾感。切叩168。

厅、县一般会核准地方头人的禀举，并要求被禀举人要签具任充状，禀举的地方头人也要签具保结状。但有时官府对地方头人的禀举不予核准，同治十年（1871)

\vfill
\begin{BookNotes}
\BookNote{168.}{「淡新档案」，第12217-2号。}
\end{BookNotes}

% 原图说　图4　光绪二年九芎林等庄绅衿、垦户、乡耆、铺户等，禀举姜荣椿为总理。

% 原书第 478 页（PDF 第 491 页）
\BookPage{478}{491}
\noindent 十一月竹南二保铜锣湾等庄生员赖廷芳、邱龙章、吴文澜、陈德熙、李宴林，监生刘蓉光、刘耀先、彭嘉谟、吴燿南、刘森元，暨铺、佃户等向淡水同知禀举李逢年为铜锣湾等联庄公局局首，淡水同知周式濂以恐生流弊，反生害民，驳斥不准。169同年十二月改以保铜锣湾等庄总理彭继生、绅衿等禀举李逢年为联庄约首，请发给谕戳，但淡水同知周式濂又以「恐所举非人，则假公济私，武断乡曲，百姓反遭鱼肉」170驳回。同治十一年（1872）四月铜锣湾等庄总理彭继生、隘首金树福、廪生刘廷珍暨垦户、绅耆等又再举李逢年为铜锣湾总理，周式濂还是没有核准，这种三举三驳的情形不多见，但也不是罕例。为蒙混一乡职，如此努力钻营，口说为地方公益，难怪周式濂怀疑其居心。171

(2)保结状:保结状是一种具结，是在保举某人为某总理（如图5)

\begin{BookQuote}
保结的内容有三部分，一为担保被举人堪充，具备适格的条件，所担保是「公事谙练，处事和睦」、172「诚实秉正，公事熟悉」、173「为人正直」、174「素行公平和众」、175「众皆悦服」。176二是担保被举人的行为与不行为，就担保其行为而言「勤慎奉公」、177「守公奉公」、178「谨慎办公」、179「小心奉公」、180「实力奉公」。181就担保其不行为而言「不
\end{BookQuote}

\vfill
\begin{BookNotes}
\BookNote{169.}{「淡新档案」，第12212-4号。}
\BookNote{170.}{「淡新档案」，第12212-1号。}
\BookNote{171.}{「淡新档案」，第12212-5号。}
\BookNote{172.}{「淡新档案」，第12243-2号。}
\BookNote{173.}{「淡新档案」，第12235-2号。}
\BookNote{174.}{「淡新档案」，第12231-2号。}
\BookNote{175.}{「淡新档案」，第12222-3号。}
\BookNote{176.}{「淡新档案」，第12217-3号。}
\BookNote{177.}{「淡新档案」，第12235-2号。}
\BookNote{178.}{「淡新档案」，第12234-2号。}
\BookNote{179.}{「淡新档案」，第12217-3号。}
\BookNote{180.}{「淡新档案」，第1227-2号。}
\BookNote{181.}{「淡新档案」，第12203-24号。}
\end{BookNotes}

% 原书第 479 页（PDF 第 492 页）
\BookPage{479}{492}
\BookArt{p492-01.jpg}{0.564}{0.422}

\noindent 敢玩误」、182「不敢舞弊玩误」、183「不敢妄为情弊」、184「无藉色滋端」、185「不敢怠玩旷职」、186「罔敢过误」、187「不敢违误」188。三是担保保结人本身而言「无虚捏冒保」、189「不敢冒结」、190「惟元等是问」、191「不敢擅保误公」、192「愿甘坐罪」、193「不敢冒保」。194

\vfill
\begin{BookNotes}
\BookNote{182.}{「淡新档案」，第12235-2号。}
\BookNote{183.}{「淡新档案」，第12234-2号。}
\BookNote{184.}{「淡新档案」，第12232-3号。}
\BookNote{185.}{「淡新档案」，第12222-3号。}
\BookNote{186.}{「淡新档案」，第12217-3号。}
\BookNote{187.}{「淡新档案」，第12203-7号。}
\BookNote{188.}{「淡新档案」，第12202-16号。}
\BookNote{189.}{「淡新档案」，第12243-2号。}
\BookNote{190.}{「淡新档案」，第12235-2号。}
\BookNote{191.}{「淡新档案」，第12234-2号。}
\BookNote{192.}{「淡新档案」，第12230-2号。}
\BookNote{193.}{「淡新档案」，第12229-3号。}
\BookNote{194.}{「淡新档案」，第12218-22号。}
\end{BookNotes}

% 原图说　图5　胡济清等保举蔡文光保结状

% 原书第 480 页（PDF 第 493 页）
\BookPage{480}{493}
\BookSubheading{(3)认充结状：}

\begin{BookQuote}
认充结状是被举人表示愿意充当总理的文书，应是在验充后才提出，但官府为便宜行事，每于保举的同时，先行提出认充结状。（如图6)
\end{BookQuote}

被举人呈任充结状，表示充当总理的意愿，并要宣示他充当总理后的行为与不行为。其内容也分三方面:一是就其行为而言有「守分奉公」、195「小心奉公」、196「勤慎办公」、197「实心奉公」、198「随传随到」、199「尽心赴义」、200「言缉盗匪、保固地方」、201「遵照宪章、联甲保防」。202二是就其不行为言有「不敢玩误」、203「不敢偏袒玩误」、204「不敢过误」、205「不敢疏懈贻误」、206「不敢怠玩旷职」、207「不敢玩误情弊」、208「不敢武断乡曲，作奸犯科」。209三是就其充任言有「不敢冒充」、210「愿甘坐罪」、211「甘愿治罪」、212「不敢玩误冒充」。213

认充结状是一种切结行为为，所以又称「切结状」，214而这种切结只是形式，在淡新档案中，并没有发现因违反切结，而被处罚的案例。

\vfill
\begin{BookNotes}
\BookNote{195.}{「淡新档案」，第12202-4号。}
\BookNote{196.}{「淡新档案」，第12202-9号。}
\BookNote{197.}{「淡新档案」，第12203-23号。}
\BookNote{198.}{「淡新档案」，第12209-25号。}
\BookNote{199.}{「淡新档案」，第12218-22号。}
\BookNote{200.}{「淡新档案」，第12217-4号。}
\BookNote{201.}{「淡新档案」，第12218-23号。}
\BookNote{202.}{「淡新档案」，第12222-2号。}
\BookNote{203.}{「淡新档案」，第12202-4号。}
\BookNote{204.}{「淡新档案」，第12203-6号。}
\BookNote{205.}{「淡新档案」，第12218-22号。}
\BookNote{206.}{「淡新档案」，第12214-3号。}
\BookNote{207.}{「淡新档案」，第12217-4号。}
\BookNote{208.}{「淡新档案」，第12218-23号。}
\BookNote{209.}{「淡新档案」，第12229-2号。}
\BookNote{210.}{「淡新档案」，第12214-3号。}
\BookNote{211.}{「淡新档案」，第12229-2号。}
\BookNote{212.}{「淡新档案」，第12232-2号。}
\BookNote{213.}{「淡新档案」，第12237-4号。}
\BookNote{214.}{「淡新档案」，第12236-3号。}
\end{BookNotes}

% 原书第 481 页（PDF 第 494 页）
\BookPage{481}{494}
\BookArt{p494-01.jpg}{0.764}{0.456}

2、验充：

县官在受理地方头人选举总理的禀文后，有权给予查核与批驳，但若认为被保举人适任，就会饬令差役传被保举人到衙署讯验，县官验得被保举人堪充总理，就会发给谕帖与戳记，这一程式就称为验充。而验充可分为讯验与发给谕戳两方面来说明，

\BookSubheading{（1）讯验：}

\begin{BookQuote}
县官认为被保举人适任总理，才需要传被保举人到衙署，接受县官的讯问，讯问是以公开方式进行，程式是先由被保举人说明相关年籍资料、职业、家庭状况，并表明认充意愿。接续由县官提出相关讯问，215而讯验的目的在审查被保举人是否适任总理一职，经讯验后若堪充，县官会当堂宣告，并在堂上训谕总理「守法奉公，毋得生事扰民」216或「庄内如有口舌
\end{BookQuote}

\vfill
\begin{BookNotes}
\BookNote{215.}{「淡新档案」，第12202-20号。}
\BookNote{216.}{「淡新档案，第12222-5号。}
\end{BookNotes}

% 原图说　图6　九芎林街总理姜荣春认充状

% 原书第 482 页（PDF 第 495 页）
\BookPage{482}{495}
细故，立即排难解纷，妥为调处，不得唆弄是非，借端需索」217等语，而

\begin{BookQuote}
县官也会在讯验后，提出工作指示「赶紧清庄、联庄，于扼要路径，搭盖更楼，派拨庄民支更把路；实力奉行」，218最后发给谕帖和戳记。
\end{BookQuote}

\BookSubheading{(2)、发给谕帖及戳记：}

A、谕帖：

\begin{BookQuote}
谕帖是官府委派总理的谕令，属一种人事命令，由总理收执保管，做为总
\end{BookQuote}

理执行其职务的依据，而总理退办或斥革，谕帖需要缴回。219（谕帖如图7）

\begin{BookQuote}
谕帖是总理的人事命令，没有谕帖总理在执行职务时，合法性与正当性都会遭到质疑，同治八年九芎林总理徐安邦就向淡水同知禀称：「虽给发戳记，其谕未蒙给发，凡遇庄中大小事务，俱称无谕，不能信从，办理维艰。隶未蒙给谕，所到各庄办理，庄民言无凭据，不敢向前办理。不蒙恩
\end{BookQuote}

准给谕执凭，恐装中所办大小事务难服人心」。220

B、戳记：

\begin{BookQuote}
戳记是厅县管发给总理，以表征其职务的印记，是方形木质，上方是给戳厅、县官的职衔，例如：淡水分府周。左右侧是某某街庄某某总理戳记，正中为花押，花押的文字是厅县官自订，一般具有勉励的意思，如「奉公」、「谨慎」、「清白」，也有直书「总理」字样（如图8）。
\end{BookQuote}

总理在退办或被斥革时，县官会吊销其谕戳，在县官更迭时，总理也要更换新戳记，清代台湾官府非常重视换戳，在淡新档案中记有新旧户粮税总书吕祥与李登间交接各案卷中有「饬吊各总董谕戳缴换卷壹宗」，221总理戳记的缴换有专卷管理。

总理并没有任期限制，有当四十多年的情形，222会换戳记是新任官的要求，换戳的理由是新任官关切地方，总理是地方最重要的乡职，对地方利益影响最大，所以在

\vfill
\begin{BookNotes}
\BookNote{217.}{「淡新档案」，第12214-9号。}
\BookNote{218.}{「淡新档案」，第12203-9号。}
\BookNote{219.}{「淡新档案」，第12238-1号。}
\BookNote{220.}{「淡新档案」，第12209-28号。}
\BookNote{221.}{「淡新档案」，第11306-1号。}
\BookNote{222.}{「淡新档案」，第12203-4号。}
\end{BookNotes}

% 原书第 483 页（PDF 第 496 页）
\BookPage{483}{496}
\BookArt{p496-01.jpg}{0.933}{0.366}

\BookArt{p496-02.jpg}{0.927}{0.336}

\noindent 就任之初，藉换戳汰换不适任的总理，这当然是为政者的用心，所以戳记上都有发给官员的职衔，官员更动戳记也一定要换新。但换戳也可能成为需所的手段。道光七年（1827）任宜兰通判的洪煌就以「各通、土所领戳记行用已久，字画糊混，着即缴销


% 原图说　图7　给大溪墘庄总理刘兆荣谕帖文

% 原图说　图8　乃丰顺、庄文蔚、李喜禄总理戳记

% 原书第 484 页（PDF 第 497 页）
\BookPage{484}{497}
\noindent 换给，每戳派出番银五十元，以供费用」。223通事、土目是长行戳按规定不用换戳，洪煌还是假借名目换戳，且每戳要价五十元，对通事、土目是不小的负担。清代台湾的官吏不仅洪煌如此，是一般都是如此，所以「每遇新官到任，即应换给戳记，备缴费银，致各办公之番竭棉莫措」。224清代自嘉庆十七年设治，到光绪二十一年移交日本，共在宜兰统治八十三年，历任通判和县令共六十任，平均一任不到一年四个月，如每次通判或县令上任都要换戳，通事、土目一定没办法负担费银，因当时宜兰已设番总理主理番社事务，通事、土目位低权轻，所以在道光十八年向通判陈情，希望通事、土目「毋庸官为饬充，以杜吊验换戳之弊」，以「免花费之累」，当时通判阎圻免通事、土目按任缴换戳记，为昭信守，并将历任换戳原卷注销，更在城隍庙立「严禁勒派社番换戳规费碑记」225示禁。阎圻表现出十分诚意，也立碑示禁，但换戳依然在进行，只隔十多年，在咸丰八年(1855）东、西势各社总通事畹吱什美龙等又为新任官到任通、土，就要缴换记，备贴费银的事提出陈情，通判杨承泽又立「严禁书差勒派社番规费碑记」，226重申「不得巧立吊验名色，又开需索之端」，但无奈法立弊生，立碑示禁，可能有短暂效果，但清代台湾换戳需索还是非常盛行，而总理是地方最重要的乡职，依规定又要换戳，所以换戳的规费应当是五十元的倍数。

\BookSubheading{(三）清代台湾乡职有严格的退革机制}

清代台湾总理被官府视为「在官人役」，227所以官府有一套管理程式，进用要经禀举、保结、认充、验充、给戳、给谕等任命程式，总理的离职也有离职的程式，不是来去自由。

1、退革的分类：

总理的离职称退革，而退革分为退职与斥革，由总理主动请辞就称退办，退办都

\vfill
\begin{BookNotes}
\BookNote{223.}{何培夫主编，〈严禁勒派社番换戳规费碑记〉《台湾地区现存碑碣图志一宜兰县·基隆市篇》（台北：国立中央图书馆台湾分馆，1999年），页26。}
\BookNote{224.}{同上注，页26。}
\BookNote{225.}{同上注，页25-27。}
\BookNote{226.}{同上注，页6-8。}
\BookNote{227.}{「淡新档案」，第12206-1号。}
\end{BookNotes}

% 原书第 485 页（PDF 第 498 页）
\BookPage{485}{498}
\noindent 是出于主观因素，而被动被除去总理职务为斥革，斥革绝大多数出于外在原因。

\BookSubheading{(1）总理退办：}

\noindent 总理因有不能再从事公务的理由，主动提出辞呈，经厅、县官核准就是所谓退办，退办是总理主动提出，厅、县官被动核准。总理退办在淡新档案中并不多见，而有就是健康因素，唯一例外的案例，是宝斗仁总理黄文绣，因变名冒充被地方另一派系举发，几经缠讼，最后经协商，由黄文绣主动请辞，以退办收场，将总理一职让给反对派。228在档案中反倒是总理在任上亡故的案例较多，这显示总理不轻易退办，除非逼不得已。而总理无论以任何私人理由请辞，厅、县官都核可退办没有慰留，但退办总理会举他人或子继任，厅、县官大都也会探纳。

\BookSubheading{(2）总理斥革：}

\noindent 总理的斥革，是厅、县官认为总理不适任，主动革去他的职务。也有地方头人联名要求厅、县官革去总理的职务。229斥革是总理被动被解除职务，退办是总理主动要求解除职务，所以退办较理性平和没有纠纷，斥革非出于总理自愿，总理为保住地位，会有救济行为，由地方头人联名具禀，希望依旧充当总理，230但厅、县官都不会准他复充。只是被斥革的总理不会就此灰心，稍有机会，随时都想复出，所以有「承担是任四十余年，三退三举，无喜无愠，不伎不求」231的情形出现。总理的斥革也常成为地方派系的争夺斗争，光绪七年（1881）吞霄街总理张凤岐被斥革，由林愈熏接充，到光绪九年（1883）拥护张凤岐的生员曾肇祯及监生、武生等五人联名指控林愈熏「胆恃威权在握，曲直在己，致使雀角鼠牙争，既不能化大成小，而户婚命盗控，反多以云集鳞聚」，232要求斥革林愈熏，另举妥人接充。随后张凤岐这派又由街庄正与铺户再联名，要求斥革林愈熏，

\vfill
\begin{BookNotes}
\BookNote{228.}{「淡新档案，第12218-20号。}
\BookNote{229.}{「淡新档案」，第12224-1号。}
\BookNote{230.}{「淡新档案」，第12203-4号。}
\BookNote{231.}{「淡新档案」，第12203-4号。}
\BookNote{232.}{「淡新档案」，第12224-4号。}
\end{BookNotes}

% 原书第 486 页（PDF 第 499 页）
\BookPage{486}{499}
\noindent 让张凤岐复充，并指责林愈熏对吞霄街总理「早欲谋充已久，乘机布摆接办」，而且「自承充以来，设馆收词，遇事生波」，还细数林愈熏「额外勒派」、「借此妄禀」、「出单勒派」、「牵牛抢物」、「任听匪徒出没」等罪行。233林愈熏也有一派地方生员职员铺户出面为他辩解，联名要求总理仍给林愈熏办理。234地方为争夺总理分派互斗，在斥革与复充间选边站。

2、退革的程序：

\BookSubheading{(1)退办的程序：}

\noindent 退办的程式较为单纯，由总理提出退办申请，厅、县官批准后，令他选举妥人接充，并缴销戳记，235就算完成退办程序。也有总理先遴选出接充人选，再提出退办申请，如此厅、县官会要求先缴销谕戳，接充人选另案验充。236

\BookSubheading{(2)斥革程序：}

\noindent 总理的斥革，有由地方头人联名禀请斥革，237也有由差役禀请斥革，238最多的是由厅、县官主动提出斥革，同治七年（1868）淡水同知严金清一次就斥革九芎林等庄总理姜荣茂、刘维兰、徐振纲，树杞林庄总理彭捷科，六张犁庄总理林炳华等五人。239总理的斥革由厅、县官示革后，命差役吊销谕戳，并出示晓谕各庄，宣示总理已经「示革在案，嗣后凡有地方公事」，不得仍向原总理「私相受授，串同索扰，致干并究」，240总理的斥革程式大致如此。官府会重视

\vfill
\begin{BookNotes}
\BookNote{233.}{「淡新档案」，第12224-1号。}
\BookNote{234.}{「淡新档案」，第12224-3号。}
\BookNote{235.}{「淡新档案」，第12202-13号。}
\BookNote{236.}{「淡新档案」，第12214-1号。}
\BookNote{237.}{「淡新档案」，第12225-1号。}
\BookNote{238.}{「淡新档案」，第12229-4号。}
\BookNote{239.}{「淡新档案」，第12209-13、20号。}
\BookNote{240.}{「淡新档案」，第12209-13号。}
\end{BookNotes}

% 原书第 487 页（PDF 第 500 页）
\BookPage{487}{500}
\noindent 退革总理谕戳的缴销，甚至晓谕各庄，是因总理出身地方，又久任其职，在地方有错综复杂的关系与影响力，也有案例显示总理被斥革，但谕戳没有吊销，被斥革的总理依然在地方执行职务，是「既格其名，而未除其实」，241所以吊戳、晓谕，是为不让被斥革的总理有再介入地方事务的机会。

3、退革的理由：

\BookSubheading{（1）退办理由：}

\noindent 退办的事由，以健康因素为主，道光二十三年（1843）竹堑本城西门总理洪德梁，是以「废疾不能办事」242退办。道光二十四年（1844）竹堑本城北门总理郭尚茂，也是以「患病不能办事」243退办。同治十一年中港头份等庄总理黄玉堂，以「年迈六旬」244退办。光绪五年（1879）宝斗仁等庄总理黄文绣，因「变名冒充」被地方不同派系举发，自请退办，这是退办的鲜例。光绪八年蛤子市庄总理张焕彩，以「年七旬有余，难以办理」245退办。在案例中总理的退办，几乎是以老病不能办事为主要理由。

\BookSubheading{(2)斥革理由：}

\noindent 斥革的事由比较复杂而多样，有「误公」、「需索」、「妄为」、「滋扰」、最严重的是「私受投词」，也有不需要理由就可斥革，现以实例说明：

A、误公：道光二十二年（1842）吞霄总理郑妈观，因「误公」246被斥革。道光二十三年宛里街总理郑文博，因「办公不力」被斥革。247光绪十二年

\vfill
\begin{BookNotes}
\BookNote{241.}{「淡新档案」，第12225-1号。}
\BookNote{242.}{「淡新档案」，第12202-13号。}
\BookNote{243.}{「淡新档案」，第12202-18号。}
\BookNote{244.}{「淡新档案」，第12214-1号。}
\BookNote{245.}{「淡新档案」，第12236-4号。}
\BookNote{246.}{「淡新档案」，第12203-4号。}
\BookNote{247.}{「淡新档案」，第12203-10号。}
\end{BookNotes}

% 原书第 488 页（PDF 第 501 页）
\BookPage{488}{501}
\noindent （1886）后垅街等庄总理杜和安，因经「三款谕送，概不办理，查造丁口清册」248被斥革。光绪十四年（1888）接充林思赞的铜锣湾等庄总理曾道隆，也是「奉公疲玩」249被斥革。在淡新档案中，总理被斥革，以「误公」为事由所占比例最多。

B、需索：同治七年（1868）九芎林等庄总理姜荣茂、刘维兰、徐振纲，树杞林庄总理彭捷科，六张犁庄总理林炳华，中坜总理黄国安等六人，因「遇事需索，唯利是图」250被斥革。光绪十一年（1885）铜锣湾等庄总理林思赞，因「鱼肉乡民」251被斥革。

C、妄为：光绪十年（1884）大溪墘庄总理钟廷英，因「遇事招摇，假公犯法」。252光绪十五年（1889）九芎林等庄总理曾克明，因「叠次妄为」253被斥革。

D、滋扰：光绪九年（1883）九芎林等庄总理林换荣，因「遇事滋扰」254被斥革。

E、私受投词：设馆收词是侵犯到厅、县官的权威，一向被列为厉禁，在总理的验充堂谕或谕帖中也一再申明，但私受投词的情形还是屡有发生，光绪九年(1883）吞霄街总理林愈熏，因「私受投词，借端骚扰」255被斥革。

F、通饬裁革：光绪九年（1883）台北知府陈星聚通令地方裁革总理，「府宪通饬裁革，

\vfill
\begin{BookNotes}
\BookNote{248.}{「淡新档案」，第12229-4号。}
\BookNote{249.}{「淡新档案」，第12238-1号。}
\BookNote{250.}{「淡新档案」，第12209-13、20、21号。}
\BookNote{251.}{「淡新档案」，第12227-1号。}
\BookNote{252.}{「淡新档案」，第12225-1号。}
\BookNote{253.}{「淡新档案」，第12242-1号。}
\BookNote{254.}{「淡新档案」，第12223-1号。}
\BookNote{255.}{「淡新档案」，第12224-1号。}
\end{BookNotes}

% 原书第 489 页（PDF 第 502 页）
\BookPage{489}{502}
\begin{BookQuote}
即无事故，亦可革」，256所以纵使没有事由，厅、县官也可斥革总理。
\end{BookQuote}

五结语

\begin{BookQuote}
「台湾为荒服之地，我先民之来居聚者，耕渔并耦，无诈无虞，出入相友，守望
\end{BookQuote}

相助，疾病相扶持，但有乡约，而无国法，固不知其几何世也。」257到清初台湾还处在拓垦时期，官方治理的脚步远落后于垦民拓垦速度，所以「官未开而民先开」。258在此状况下垦民自成聚落，也进行自我管理，长期实行没有官府的乡治。这在康雍干时期还普遍存在，而嘉庆、道光以后全台几已垦透，厅、县也已普遍设立，「但有乡约，而无国法」的地方治理模式已少之又少。

在康雍干时期，朝廷逐渐在台湾设立厅、县，只是台湾远隔大洋，清朝的治理态度又消极，再说台湾当时户口不多但辖区辽阔，如「诸罗建自本朝，民户二千四百三十有六；方之前代，其在中、下之间乎·然论其地，则战国时之宋、卫、中山弗及也·自南以北，袤千有余里，往返动以月计；山深水阻，宵小易潜·县令一耳目之力，势不得不有所寄，无能出胥役、保长、通事者矣」。25因此「台属有司，率多不亲理民事，皆批乡约、保长、管事、书役查覆」实属必然，也促发清代台湾乡治的蓬勃。

清代台湾有三层乡治组织，分别是坊里保与图、庄及牌、甲。并置头人总理、董事、街庄正、地保、族长负责管理。而乡职有非常重的职权「凡人民的迁徙、职业、婚嫁、生死，均要报于总理」。且「台属有司，率多不亲理民事，皆批乡约、保长、管事、书役查覆」，所以有「总理不职，地方攸关」260的说法。

法制史学者戴炎辉也在他的代表作《清代台湾之乡治》中如是说:「总理系官最重视的乡职，称为一堡之正；绅耆等于保举总理之禀内，则称之为街庄之首，总理

\vfill
\begin{BookNotes}
\BookNote{256.}{「淡新档案」，第12224-3号。}
\BookNote{257.}{连横，《台湾通史》（南投:台湾省文献委员会，1992年)，页283。}
\BookNote{258.}{《福建台湾奏折》（南投:台湾省文献委员会，1997年），页1。}
\BookNote{259.}{周钟瑄，《诸罗县志》（南投：台湾省文献委员会，1993年），页51。}
\BookNote{260.}{「淡新档案」，第12224-1号。}
\end{BookNotes}

% 原书第 490 页（PDF 第 503 页）
\BookPage{490}{503}
\noindent 系境内之核心的存在，其职务则该括全面，其被称为总理，职此之故」；又说「总理及地保系乡治上最重要的职务，总理本为地方自治团体的首席，且为其执行人」。261总理是清代台湾最重要又最被重视的乡职，是「街庄之首」更是「地方自治团体的首席」，是地域内核心的所在，且其职务是眩括全面。

「总理不职，地方攸关」且其职务赔括全面，说明清代台湾乡职的重要性。而后总理乡职发展成普设、常设制度。官府造有各保总理、董事清册，或有出缺，要立即递补。乡职已成行政一环，清代台湾地方政务几乎委任乡职在推行。

\vfill
\begin{BookNotes}
\BookNote{261.}{戴炎辉，《清代台湾之乡治》（台北:联经出版事业公司，1979年）页21-22。}
\end{BookNotes}
